\input{glyphtounicode}
\pdfgentounicode=1
% =============================================================================
% Globale Normierung in Joshis Arithmetischer Teichmüller-Theorie
% Einordnungs-/Klassifikationspapier -- Reparaturfassung v1.2 (Deutsch)
% Projekt: DRAFT__Global_Normalization_Contract_Lemma  (.RESEARCH/.LAB/.HCT)
% Deutsche Fassung von GNC_Assessment_v1_en.tex (1:1-Sektions- und
% Absatzparität; wörtliche Joshi-Zitate bleiben englisch).
% Basis: OUTLINE.md; _proof-notes/LEMMA_FORMULIERUNG_V1_2026-08-12.md;
%        P1_QUELLENPRUEFUNG_NORMIERUNG_2026-08-12.md;
%        P2_TERM_KLASSIFIKATION_THM6101_2026-08-13.md;
%        R-E2_MASSPINNUNG_2026-08-13.md
% v1.1 (13.08.2026): zweigkonditionale Revision. Jedes Verdikt, das von der
%        Auswahlregel für lambda_y abhängt, ist jetzt für den Zweig markiert,
%        unter dem es gilt (Trilemma, Abschnitt 3.5); Lemma 1 (Starrheit des
%        Produktformel-Multiplikators) in Abschnitt 3.6 ergänzt; Abschnitt 7 neu
%        geschrieben (der frühere Dreier-Lesarten-Rahmen ist überholt);
%        KI-Offenlegung auf den Pipeline-Standard per \input umgestellt.
%        Bericht: _results/V11_TRILEMMA_REVISION_2026-08-13.md
% Dies ist KEIN Beweispapier. Siehe den Geltungsbereichs-Kasten nach dem Abstract.
% =============================================================================

\documentclass[11pt,a4paper]{article}

\usepackage{cmap}
\usepackage[utf8]{inputenc}
\usepackage[T1]{fontenc}
\usepackage{lmodern}
\usepackage{microtype}
\usepackage{xurl}
\usepackage[ngerman]{babel}
\usepackage[a4paper,margin=2.7cm]{geometry}
\usepackage{amsmath}
\usepackage{amsthm}
\usepackage{amssymb}
\usepackage{array}
\usepackage{booktabs}
\usepackage[colorlinks=true,linkcolor=black,citecolor=black,urlcolor=blue,unicode=true]{hyperref}
\def\HyperDestNameFilter{ger.}


\babelhyphenation[ngerman]{
  Ara-ke-lov Arith-me-ti-coid Arith-me-ti-coids Arith-me-ti-coi-de
  Arith-me-ti-coi-den Ho-lo-mor-pho-id Ho-lo-mor-pho-ide Ho-lo-mor-pho-iden
  Nor-mie-rung Nor-mie-rungs-ko-or-di-na-te Pro-dukt-for-mel-Hy-per-ebe-ne
  Be-wer-tungs-ska-la Be-wer-tungs-ska-len Zahl-kör-per-ele-men-te
  Kon-struk-ti-on Men-gen-sta-bi-li-tät
}

% v1.1: Die Revision hat mehrere Absätze mit dicht gesetzten, unteilbaren
% Mathe-Atomen ergänzt; etwas zusätzliche Dehnung hält sie im Satzspiegel.
\emergencystretch=1.5em

% ragged-right Absatzspalten für die beiden Klassifikationstabellen
\newcolumntype{P}[1]{>{\raggedright\arraybackslash}p{#1}}

\hypersetup{
  unicode=true,
  pdfencoding=unicode,
  psdextra,
  breaklinks=true,
  pdftitle={Globale Normierung in Joshis Arithmetischer Teichmüller-Theorie},
  pdfauthor={Lukas Geiger},
  pdfsubject={Einordnung der Normierungs-Identifikation in ATS IV Theorem 6.10.1},
  pdfkeywords={Arithmetische Teichmüller-Theorie, Globale Normierung, Produktformel, Artin-Whaples, Log-Volumen, Frobenius-Verschiebung}
}

% --- theoremartige Umgebungen -------------------------------------------------
\theoremstyle{definition}
\newtheorem{question}{Frage}
\theoremstyle{plain}
\newtheorem*{derivationD}{Konditionale Herleitung D-1}
\newtheorem{lemmaR}{Lemma}
\theoremstyle{definition}
\newtheorem*{findingF}{Befund F-NEU-1}
\newtheorem*{criterionK}{Kriterium K (Kandidat)}

% --- Notation -----------------------------------------------------------------
\newcommand{\arith}{\operatorname{arith}}
\newcommand{\hol}{\operatorname{hol}}
\newcommand{\ordd}{\operatorname{ord}}
\newcommand{\Vol}{\operatorname{Vol}}
\newcommand{\LogVol}{\operatorname{LogVol}}
\newcommand{\Thet}{\widetilde{\Theta}}
\newcommand{\Ynum}{Y_L}

% Jede Quellenangabe ist an ein benanntes Extrakt gebunden, weil die vier
% Extrakte unabhängige Zeilennummerierungen tragen.
\newcommand{\anchor}[2]{(#1, Extrakt Z.~#2)}

\title{Globale Normierung in Joshis Arithmetischer Teichmüller-Theorie:\\
Was die Identifikation trägt --- und was das Theorem trägt}
\author{Lukas Geiger\\[2pt]
\small Unabhängiger Forscher, Bernau, Deutschland\\
\small ORCID: 0009-0005-7296-1534}
\date{2. Oktober 2026 --- Reparaturfassung v1.2}

\begin{document}
\maketitle

\begin{abstract}
\noindent
Diese Einordnung fragt, was dieselbe globale Normierung in
Theorem~6.10.1 von Joshis ATS~IV (arXiv:2403.10430v2) trägt. Unter der ausdrücklich
stellenweisen Artin--Whaples-Lesart werden Zahlkörperbeträge definitorisch
identifiziert; die Auswahlregel für die Koordinate $\lambda_y$ bleibt eine offene
Korrespondenzfrage. Ein elementares Produktformel-Lemma beschränkt Multiplikatoren
der Standardbewertungen auf einen globalen Skalar. Der Ausschluss einer festen
Produktformel-Hyperebene bleibt an die Annahmen (F1)--(F2) gebunden.
Die drei Normierungszweige dienen der Diagnose und sind nicht erschöpfend.
Diese Reparaturfassung trennt die quellenbehauptete Frobenius-Stabilität einer
kollationierten Familie vom ungeprüften Transport zu den beiden indizierten
Hüllen in E2. Sie berichtigt drei Quellenlesungen: ATS~IV nennt
$\overline L^*$ als Höhendomäne, E2 setzt Beträge von Log-Volumina gleich,
und der algebraische holomorphe Hull aus ATS~III ist allgemein kein konvexer
Abschluss. Ein expliziter Gegenfall ergibt
$\operatorname{conv}(\mathbb Z_p\cdot1)=\mathbb Z_p\cdot1
\subsetneq\mathcal O_E=\operatorname{hull}(\mathbb Z_p\cdot1)$ in einer
unverzweigten quadratischen Erweiterung $E/\mathbb Q_p$. Das widerlegt weder
die speziellen Aussagen über Tensorregionen noch die ATS-Schranken.
Normiertes additives Haar-Maß bleibt von Potenzen eines Absolutbetrags getrennt:
$m(p\mathbb Z_p)=p^{-1}$ wird durch einen geänderten Bewertungsexponenten
nicht zu $p^{-c}$. Die frühere positive Behauptung einer frameübergreifenden
Maßfestlegung wird deshalb in jedem Zweig zurückgenommen.
Ein indizierter, hull- und maßverträglicher Transport wird als mögliche
hinreichende Route benannt, die stärker als die angeschriebene Betragsgleichheit
ist; er ist weder bewiesen noch als notwendig ausgewiesen. Auswahlregel,
indizierter E2-Transport, verträgliche Maß- und Gewichtsbindungen sowie
Ordnungsvergleich bleiben offen. Ein Beweis oder eine Widerlegung von $abc$,
Mochizukis Korollar~3.12 oder Joshis Abschätzungen wird nicht beansprucht.
\end{abstract}

\clearpage
% -----------------------------------------------------------------------------
% GELTUNGSBEREICH UND BEHAUPTUNGSGRENZE
% -----------------------------------------------------------------------------
\begin{center}
\fbox{%
\parbox{\dimexpr\textwidth-2\fboxsep-2\fboxrule\relax}{%
\medskip
\textbf{Geltungsbereich und Behauptungsgrenze.}\par\medskip
\begin{itemize}\setlength{\itemsep}{2pt}
\item Diese Arbeit beweist \emph{nichts} über die $abc$-Vermutung, über die
  inter-universelle Teichmüller-Theorie, über Mochizukis Korollar~3.12 oder über
  die Korrektheit von Joshis Schranken. Sie enthält weder einen Beweis noch eine
  Widerlegung von irgendetwas davon.
\item Sie \emph{verifiziert keine Schranke}. Keine Beweiszeile von ATS~IV
  Theorem~6.10.1, von Korollar~9.11.1.1 oder von Korollar~3.12 wurde Schritt für
  Schritt nachvollzogen; keine Abschätzung wurde nachgerechnet.
\item Was sie tut: Sie \emph{klassifiziert Definitionen und die Struktur der
  Begründung} in einem benannten Theorem und gibt für jede seiner Komponenten an,
  ob diese definitorisch, durch die Konstruktion geliefert, zweigkonditional oder
  offen ist.
\item Sie bezieht keine Position in der Auseinandersetzung zwischen Joshi und
  Scholze--Stix~\cite{ScholzeStix2018} über die Beweglichkeit der
  Produktformel-Hyperebene. Wo diese Auseinandersetzung berührt wird, wird sie
  berichtet, nicht entschieden.
\item Die in Abschnitt~\ref{sec:b2outlook} skizzierte strukturelle Analogie ist
  eine Arbeitshypothese des vorliegenden Programms über die \emph{Gestalt von
  Argumenten}. Sie ist keine mathematische Verbindung und insbesondere keine
  Aussage über die Riemannsche Vermutung.
\item Der zentrale offene Punkt dieser Fassung --- dass die von uns gelesenen
  Passagen keine Regel nennen, die die Normierungskoordinate eines bewegten
  Arithmeticoids auswählt (Abschnitt~\ref{sec:trilemma}) --- ist eine
  \emph{Definitionslücke}: eine Regel, die wir nicht gefunden haben und die
  durchaus in nicht gelesenem Material stehen kann. Sie ist keine Widerlegung von
  ATS~III Korollar 4.2.2.4 und keinerlei Fehlerbehauptung. Lemma~1
  (Abschnitt~\ref{sec:rigidity}) schränkt ein, welche Regeln die Produktformel
  \emph{allein} liefern kann; es behauptet nicht, dass es keine solche Regel gibt.
\item Wo eine Spannung zwischen zwei Passagen einer Quelle festgehalten wird, ist
  dies eine Einladung, sie aufzulösen, und ein Kandidat für die Korrespondenz mit
  dem Autor --- keine Behauptung, sie sei unauflösbar, und keine
  Fehlerbehauptung.
\item \textbf{Textgrundlage} sind \texttt{pdftotext}-Extrakte der arXiv-PDFs.
  Alle Zeilenangaben (\glqq Extrakt Z.~NNNN\grqq) sind Auffindhilfen in diesen
  Extrakten und \emph{keine} Paginierung der veröffentlichten Dokumente. Die
  direkte visuelle Rückprüfung umfasste die Quell-PDF-Seiten zu Gleichung (5.3.3),
  Theorem~5.10.1(3)--(4), Theorem~10.20.1(3),
  Definition~4.4.2/Formel (4.4.5) und Theorem~6.10.1. Eine pauschale
  OCR-Robustheit wird nicht behauptet; insbesondere wurden die E1-Betragsstriche
  und die Richtung des Faserexponenten in den Quell-PDFs geprüft.
\end{itemize}
\medskip
}}
\end{center}

\clearpage
\tableofcontents
\clearpage

% =============================================================================
\section{Einleitung: der strittige Punkt}
% =============================================================================

\subsection{Der Befund, der zu dieser Arbeit führte}

Die vorliegende Arbeit ist der vierte Teil eines kleinen Forschungsprogramms,
dessen dritter Teil ein strukturierter \glqq Ledger-Test\grqq{}
(Anforderungsbuch-Prüfung) der Argumentationsarchitektur von ATS~IV war: eine
Prüfung Anforderung für Anforderung, welche tragenden Schritte der Arbeit
geliefert, welche referenziert und welche behauptet werden. Dieser Test gab eine
Anforderung --- diejenige zur Vergleichbarkeit der beiden in Theorem 6.10.1
eingehenden Arithmeticoide --- als nur \emph{teilweise} erfüllt zurück. Der Grund
war keine Lücke in einer Rechnung. Er war, dass die Identifikation, welche die
mittlere Gleichheit des Theorems trägt, nach dem eigenen Wort des Textes eine
Tautologie ist \anchor{ATS~IV}{3272--3273}, während die Bemerkung desselben
Theorems genau jenen stellenweisen Vergleich verbietet, mit dem man sie sonst
prüfen könnte.

Die beiden Textpole sind diese. Einerseits eröffnet das Theorem mit einer Wahl:

\begin{quote}\small
``Choose the same global normalization for the number field provided by the
arithmeticoids $y_0'$ and $y_0$. With this choice of simultaneous global
normalization one has \ldots'' \hfill\anchor{ATS~IV}{3185--3187}
\end{quote}

vorbereitet in der Einleitung durch

\begin{quote}\small
``Each arithmeticoid $y_0' = \phi(y_0)$, $y_0$ provides an isomorphic copy of our
number field and hence one has chosen the same global normalization for both the
copies, the [sic] it does not matter which arithmeticoid is used for global
computations.'' \hfill\anchor{ATS~IV}{740--743}
\end{quote}

Andererseits stellt derselbe Text unmittelbar unter dem Theorem fest:

\begin{quote}\small
``as arithmeticoids $y_0' \neq y_0$ as they differ by Frobenius at each prime
$v \in V_L$ and this, together with choice of concurrent global normalization for
the copy of the (same) number field they provide \emph{precludes direct comparison
of various local quantities}.'' \hfill\anchor{ATS~IV}{3207--3211}
\end{quote}

Beide Passagen stammen von Joshi. Sie motivieren eine Frage nach der
Identifikation der Elementskalen des Zahlkörpers, legen aber nicht die gesamte
E2-Gleichheit auf diese Identifikation. Indizierte Mengen, algebraische
Hüllen, Maße und Gewichte benötigen eigene Bindungen.

\begin{question}\label{q:main}
\emph{Behauptet} die Wendung \glqq dieselbe globale Normierung\grqq{}, wie sie in
ATS~IV Theorem 6.10.1 verwendet wird, eine Isometrie zwischen den beiden
Bewertungsskalen, oder \emph{stellt} sie eine solche per Definition \emph{her}?
\end{question}

Die beiden Antworten haben sehr verschiedene Konsequenzen. Behauptet die Wendung
eine Isometrie, dann wird etwas behauptet, dessen Beweis man zu sehen berechtigt
ist, und der Arbeitsname des Programms für die fehlende Aussage --- ein
\emph{globales Normierungs-Vertragslemma} (global normalization contract lemma)
--- wäre gerechtfertigt. Stellt die Wendung die Identität definitorisch her, dann
ist für diese begrenzte Identität der Elementskalen kein gesondertes Lemma
nötig. Was die indizierte Hull- und Volumengleichheit trägt, ist eine weitere
Frage.

Diese Arbeit beantwortet Frage~\ref{q:main} im zweiten Sinne --- konditional auf
eine Lesart, welche die Quellen offen lassen, wie Abschnitt~\ref{sec:trilemma}
darlegt --- und verfolgt anschließend die Anschlussfrage, denn in der
Anschlussfrage lag, wie sich herausstellte, die Substanz.

\subsection{Quellen, Kürzel und Zitierdisziplin}
\label{sec:sources}

Vier Dokumente dienen als Primärtext. Durchgehend verwenden wir die Kürzel der
linken Spalte; jede Quellenangabe dieser Arbeit trägt ein solches Kürzel, weil die
vier Extrakte unabhängige Zeilennummerierungen haben und eine bloße Zeilennummer
mehrdeutig wäre. Tabelle~\ref{tab:sources} fasst das Primärquellen-Korpus und die zugewiesenen Rollen zusammen.

\begin{table}[htbp]
\centering
\caption{Primärquellen-Korpus und zugewiesene Kürzel in dieser Untersuchung.}
\label{tab:sources}
\small
\begin{tabular}{@{}lll@{}}
\toprule
Kürzel & Quelle & Rolle in dieser Arbeit \\
\midrule
ATS~II & arXiv:2303.01662v3~\cite{JoshiATSII}
  & Verhalten der Bewertungen entlang der Frobenius-Faser \\
ATS~II(1/2) & arXiv:2305.10398v12~\cite{JoshiATSIIhalf}
  & Definition von $\arith(L)^{\mathrm{nor}}$; Periodenabbildung; Höhen \\
ATS~III & arXiv:2401.13508~\cite{JoshiATSIII}
  & Normierungskonvention; Hyperebene; Volumentheorie \\
ATS~IV & arXiv:2403.10430~\cite{JoshiATSIV}
  & Theorem 6.10.1 und der es umgebende Apparat \\
\bottomrule
\end{tabular}
\end{table}

In Joshis eigenen Literaturverzeichnissen erscheinen diese als \glqq [Joshi,
2023b]\grqq{} (ATS~II), \glqq [Joshi, 2023a]\grqq{} (ATS~II(1/2)) und \glqq [Joshi,
2024c]\grqq{} (ATS~III); die Zuordnung der ersten beiden wurde gegen das
Literaturverzeichnis von ATS~IV geprüft \anchor{ATS~IV}{3807--3810}, die dritte ist
die in den zugrunde liegenden Arbeitsnotizen durchgehend verwendete Identifikation.
Wo wir eine Passage wiedergeben, ist der Wortlaut wörtlich aus dem Extrakt
übernommen; mathematische Symbole wurden gesetzt, und Auslassungen sind markiert.
Zwei Lesarten auf OCR-Ebene wurden ausdrücklich getroffen: Die Zeichenfolge
\glqq ptvh-power\grqq{} in ATS~II(1/2), Bemerkung 5.9.2, wird als \glqq $p_v$-th
power\grqq{} gelesen, und die Platzierungen von Tief- und Hochstellungen in den
abgesetzten Normierungsgleichungen wurden gegen die umgebende Prosa geprüft.

\subsection{Was folgt}

Abschnitt~\ref{sec:equivariance} schärft die Frage: Die naive Form \glqq Ist $F$
invariant?\grqq{} verfehlt die Struktur, weil sich sowohl das ausgewertete Objekt
als auch die Skala bewegen, in der es ausgewertet wird.
Abschnitt~\ref{sec:definition} klärt die Bedeutung der Wendung unter der
stellenweisen Lesart der Definition von $\arith(L)^{\mathrm{nor}}$; sein
Abschnitt~\ref{sec:trilemma}
formuliert anschließend das Trilemma, das die Definition offen lässt, und
Abschnitt~\ref{sec:rigidity} beweist das Lemma, das es einschränkt.
Abschnitt~\ref{sec:D1} zeigt durch eine
elementare Herleitung unter zwei mit ihr angegebenen Annahmen, dass sich die
Produktformel-Hyperebene wirklich bewegt, was --- unter diesen Annahmen ---
die eine Lesart ausschließt, unter der die Wendung nichttrivialen Gehalt gehabt
hätte. Abschnitt~\ref{sec:terms} klassifiziert die Terme von Theorem 6.10.1 und
benennt den Trägerkandidaten. Abschnitt~\ref{sec:measure} trennt lokale Konventionen
von den offenen frameübergreifenden Maß- und Hull-Pflichten. Abschnitt~\ref{sec:fneu} formuliert die
offene Frage am Rand --- diejenige, aus der das Trilemma hervorging.
Abschnitt~\ref{sec:assessment} zieht die Konsequenzen und benennt die Grenzen.
Verdikte, die vom Zweig abhängen, sind dort als solche markiert, wo sie
ausgesprochen werden; die Übersicht, welches Verdikt unter welchem Zweig gilt, ist
Tabelle~\ref{tab:branches}.

% =============================================================================
\section{Zuspitzung der Frage: Äquivarianz, nicht Invarianz}
\label{sec:equivariance}
% =============================================================================

\subsection{Beide Seiten bewegen sich}

Ein erster Versuch, Frage~\ref{q:main} zu formalisieren, würde fragen: Ist die
Größe $F$ invariant unter einem Wechsel des Arithmeticoids? Diese Formulierung ist
falsch, und zu sehen warum, ist der erste inhaltliche Schritt.

In der Betragsgleichheit (mit den unten in E2 angegebenen Indizes)
\[
 \bigl|\log\Vol(\operatorname{hull}(\Thet^{\varphi(y_0)}))\bigr|
 =\bigl|\log\Vol(\operatorname{hull}(\Thet^{y_0}))\bigr|
\]
bewegen sich \emph{beide} Argumente: das ausgewertete Objekt und die Skala, in der
die Auswertung geschieht. Die korrekte Form der Aussage ist daher nicht Invarianz,
sondern ein indizierter, rahmenabhängiger Vergleich. Ein isometrischer Transport
 wäre eine mögliche hinreichende Route; die Betragsgleichheit allein impliziert ihn nicht. Das ist
in beide Richtungen bedeutsam. Eine Aussage dieses Typs ist nicht automatisch
trivial --- eine Isometrie zu sein, ist eine substanzielle Eigenschaft ---, sie ist
aber auch nicht automatisch leer, wenn sie gilt, anders als eine bloße Umbenennung.
Ob die Wendung \glqq identisch normiert\grqq{} eine solche Isometrie behauptet oder
sie durch eine Identifikation erzeugt, ist genau Frage~\ref{q:main}.

Die einschlägigen strukturellen Tatsachen stellt Joshi selbst fest. Die Schemata
sind isomorph:

\begin{quote}\small
``the schemes underlying these holomorphoids are isomorphic to $X/L$''
\hfill\anchor{ATS~IV}{426--427}
\end{quote}

während die analytischen Strukturen und mit ihnen die Bewertungsskalen es nicht
sind:

\begin{quote}\small
``arithmetic (and geometry) provided by the holomorphoids $\hol(X/L)_y$ and
$\hol(X/L)_{y'}$ \ldots\ occur at intrinsically different valuation scales''
\hfill\anchor{ATS~IV}{450--453}
\end{quote}

und dieser Unterschied ist keine Unannehmlichkeit der Konstruktion, sondern ihr
Zweck:

\begin{quote}\small
``It is precisely the existence of such intrinsic differences in valuation scales
which the averaging techniques \ldots\ exploit''
\hfill\anchor{ATS~IV}{457--459}
\end{quote}

\subsection{Zwei Präzisionsstufen der Frage}

Zwei Fassungen der Frage lassen sich formulieren, von ungleicher Reife.

\paragraph{(L-weak) Die Invarianzklasse.}
Sei $N$ eine feste Wahl identischer globaler Normierung für die von $y_0$ und
$\varphi(y_0)$ bereitgestellten Kopien von $L$. Charakterisiere die Klasse
\[
  G \;:=\; \bigl\{\, F \;:\; F(y_0) = F(\varphi(y_0)) \text{ ist unter } N
  \text{ \emph{beweisbar}} \,\bigr\},
\]
wobei $F$ kovariant ausgewertet wird (Objekt und Skala gemeinsam transportiert) und
\glqq beweisbar\grqq{} ausdrücklich \emph{ausschließt}, dass die Gleichheit aus der
Definition von $N$ oder aus einer Notationskonvention folgt.

Ein Beweis für $F \in G$ hätte die Gestalt einer Kette von Gleichheiten, die
$F(\varphi(y_0))$ in $F(y_0)$ überführt und dabei nur Schritte verwendet, die
entweder über nachweislich von $\varphi$ fixierte Daten laufen oder eine
nachgewiesene Kompensation zweier Änderungen ausnutzen. Der Text liefert ein
formales Modell für die erste Art: Der Beweis von ATS~IV Proposition 4.4.4
reduziert $\log(q_M) = \log(q_{L_{\mathrm{mod}}})$ auf die Gradformel
$\sum_{w\mid v} e_{w\mid v} f_{w\mid v} = [M : L_{\mathrm{mod}}]$
\anchor{ATS~IV}{1866--1928}. Das ist ein echter Beweis mit klassischen Mitteln ---
aber er zeigt Invarianz unter \emph{Körpererweiterung}, nicht unter $\varphi$, und
ist hier daher eine Schablone, keine Antwort.

Mehrere Quellenpassagen begrenzen die vorgeschlagene Klasse. ATS~IV
berichtet steigende Tate-Parameter-Beträge unter Frobenius
\anchor{ATS~IV}{567--569}, Nichtkonstanz des normierten Tate-Idel-Grades
auf $\Ynum$ \anchor{ATS~IV}{2055--2075} und eine Höhendomäne
$\overline L^*\times\Ynum$ \anchor{ATS~IV}{550--559}.
ATS~III berichtet Bewegung der Produktformel-Hyperebene
\anchor{ATS~III}{518--522}.
Das sind quellenberichtete Beispiele und Indizien, nicht vier unabhängig
bewiesene Fälle fehlender Frobenius-Invarianz. Ein angegebener
Definitionsbereich beweist weder Nichtkonstanz noch Variation entlang
einer bestimmten Frobeniusbahn. Höhenaussage und F-NEU-1 behalten die
getrennten Domänen- und Beweispflichten aus Abschnitt~\ref{sec:fneu};
D-1 behält (F1)--(F2).

\paragraph{(L-strong) Ordnungsverträglichkeit nach Summation.}
Gib ein Kriterium an, unter dem globale Invarianz \emph{Ordnungsverträglichkeit}
zwischen zwei Familien lokaler Beiträge liefert, obwohl stellenweise
Vergleichbarkeit ausgeschlossen ist. Die Nebenbedingung ist nicht erfunden; sie ist
Joshis eigene:

\begin{quote}\small
``one cannot claim that $z_p \le z_p'$ holds for all primes $p$ without violating
the middle equality of (1.5.1)'' \hfill\anchor{ATS~IV}{817--819}
\end{quote}

Diese zweite Fassung ist die weniger entwickelte der beiden: Sie benennt die
Struktur, bietet aber keine Kandidatenklasse zulässiger lokaler Verteilungen. Wir
halten das als den größten inhaltlichen Rückstand der Formulierungsphase fest; in
dieser Arbeit wird er nicht aufgelöst.

\subsection{Ein trennendes Kriterium}

Die in den Abschnitten~\ref{sec:terms} und~\ref{sec:measure} verwendete
Klassifikation ruht auf einem Kriterium, das wir als Hypothese formulieren, nicht
als Satz.

\begin{criterionK}
Eine Größe ist umso eher im tragenden Sinne invariant, je vollständiger sie über
Daten faktorisiert, die $\varphi$ nachweislich fixiert (das zugrunde liegende
Schema, Körpergrade, Mächtigkeiten der Restklassenkörper), und umso weniger, je
unmittelbarer sie die \emph{Bewertungsskala} $|\cdot|_{K_{y_v}}$ des Arithmeticoids
liest.
\end{criterionK}

Die Trennschärfe des Kriteriums lässt sich am Theorem selbst ablesen, denn
Theorem 6.10.1 enthält beide Typen nebeneinander: Seine erste Gleichheit betrifft
$q_\ell$, dessen definierende Formel über $\ordd_w(q_w)$ und $\log w$ läuft,
während seine mittlere Gleichheit $\log \Vol$ betrifft, das über
$\LogVol(\pi \mathcal{O}) = \log|\pi|$ die Skala unmittelbar liest. Diese beiden zu
trennen war das erste Ergebnis der Formulierungsphase; die Trennung durchzuführen
ist Aufgabe von Abschnitt~\ref{sec:terms}. Tabelle~\ref{tab:invariance_classes}
fasst die resultierende Taxonomie der Invarianzklassen unter einem Arithmeticoid-Shift zusammen.

\begin{table}[htbp]
\centering
\caption{Taxonomie der Invarianzklassen und des trennenden Kriteriums~K unter einem Arithmeticoid-Shift $\varphi$.}
\label{tab:invariance_classes}
\small
\setlength{\tabcolsep}{3.5pt}
\begin{tabular}{@{}P{0.07\textwidth}P{0.24\textwidth}P{0.24\textwidth}P{0.23\textwidth}P{0.14\textwidth}@{}}
\toprule
Klasse & Invarianzcharakter & Faktorisierungsbasis & Exemplare in ATS~IV & Lückenrolle \\
\midrule
\textbf{T} & Trivial / definitorisch & Faktorisiert über $\varphi$-fixierte Schema- und Graddaten ($[M:\mathbb{Q}]$, $e, f$) & $\ordd_w(q_w)$, $\log w$, $q_\ell$ (\S\,4.4) & Harmlos \\
\textbf{B} & Schranken-invariant (global) & Erhalten unter globaler Summation / Produktformel & $\sum_v \log |x|_v = 0$, $A' = CA$ & Bindend \\
\textbf{F} & Faser-äquivariant & Transformiert äquivariant entlang der Frobenius-Faser & Periodenabbildung $\operatorname{per}_{\arith}$, $\Thet$ & Strukturell \\
\textbf{N} & Nicht-invariant (Untilt) & Wertet unmittelbar die Untilt-Bewertungsskala $|\cdot|_{K_{y_v}}$ aus & Tate-Idel $\log(TI)$, Kummer-Abb. & Obstruierend \\
\bottomrule
\end{tabular}
\end{table}

% =============================================================================
\section{Die Definition unter der stellenweisen Lesart:
\texorpdfstring{$\arith(L)^{\mathrm{nor}}$}{arith(L)nor}
ist punktweise eindeutig}
\label{sec:definition}
% =============================================================================

\subsection{Drei mögliche Lesarten (nicht erschöpfend)}

Vor dem Blick in die Definition lässt die Wendung \glqq dieselbe globale
Normierung\grqq{} mindestens drei Lesarten zu. Es lohnt, diese drei zu nennen,
weil zwei aus textinternen Gründen ausgeschlossen werden können; die Liste ist
diagnostisch und beansprucht keine logische Vollständigkeit.

\begin{description}
\item[N1.] Beide Arithmeticoide sind im Sinne der stehenden Konvention normiert,
  d.\,h.\ in jedem gilt die Produktformel.
\item[N2.] Beide haben \emph{dieselbe} Produktformel-Hyperebene,
  $H_{y_0} = H_{\varphi(y_0)}$.
\item[N3.] Innerhalb der \emph{Restfreiheit}, welche die Normierungsbedingung
  offenlässt, wird eine koordinierte Wahl getroffen --- der Gedanke ist, dass
  $\prod_v |x|_{K_v} = 1$ die einzelnen $|\cdot|_{K_v}$ nicht festlegt.
\end{description}

Ein weiterer Mechanismus könnte beispielsweise die Produktformelbedingung durch
eine globale Auswahlregel ergänzen, die den verbleibenden Skalar $c$ festlegt,
ohne zur stellenweisen Artin--Whaples-Bindung von Zweig~(a)/N1$'$ zurückzukehren
und ohne die stellenselektive Freiheit von Zweig~(c)/N3 zuzulassen. Die drei
Lesarten ordnen daher die
vorliegenden Quellpassagen; sie teilen nicht alle denkbaren Auswahlregeln auf.

Lesart N1 kann für sich genommen nicht gemeint sein, denn die Konvention gilt in
der Reihe bereits universell:

\begin{quote}\small
``I will also assume throughout this paper that all arithmeticoids $\arith(L)_y$
are normalized i.e.\ $\arith(L)_y = \arith(L)^{\mathrm{nor}}_y$ \ldots\ i.e.\ the
valuations on each component of $y$ are normalized so that the Artin--Whaples
product formula (\cite{ArtinWhaples1945}) holds for each $x \in L^*$:
$\prod_{v \in V_L} |x|_{K_v} = 1$'' \hfill\anchor{ATS~III}{858--863}
\end{quote}

Unter N1 fügte der Satz in Theorem 6.10.1 nichts hinzu, was nicht ohnehin gälte.
Lesart N2 steht in sichtbarer Spannung zur Betonung der Beweglichkeit der
Hyperebene im Text selbst, auf die wir in Abschnitt~\ref{sec:D1} zurückkommen. Damit
bleibt N3 als die Lesart mit Gehalt --- und N3 ist genau diejenige, deren Definition
nicht in ATS~III oder ATS~IV steht, sondern in der früheren Arbeit.

\subsection{Was die Definition tatsächlich sagt}

Die Definition des normierten Arithmeticoids steht in ATS~II(1/2) \S\,5.3. Die
einschlägigen Passagen lauten:

\begin{quote}\small
``The fundamental global constraint \ldots\ is the validity of the product formula
for $L$ \ldots\ for the \emph{standard choice of normalizations} for valuations
$(L_v, |\cdot|_v)$ (given in [Artin and Whaples, 1945]) \ldots\ (5.3.1)
$\prod |x|_v = 1$'' \hfill\anchor{ATS~II(1/2)}{1840--1848}
\end{quote}

\begin{quote}\small
``Now suppose one is given a $y = (\psi_v) \in \Ynum$ \ldots\ Since the two
valuations $|\cdot|_v$ and $|\cdot|_{K_{\psi_v}}$ on $L_v$ are \emph{equivalent,
there exist, for each $v \in V_L$, real numbers $\lambda_v \in \mathbb{R}$} such
that (5.3.3) $(L_v, |\cdot|_v) = (L_v, |\cdot|^{\lambda_v}_{K_{\psi_v}})$.''
\hfill\anchor{ATS~II(1/2)}{1852--1858}
\end{quote}

\begin{quote}\small
``(5.3.5) $\lambda_y = (\lambda_v)$ \ldots\ \emph{is the normalization coordinate
of $y$}. When valuations induced on $\{L_v\}$ by the local data of $\arith(L)_y$
are normalized so that (5.3.4) holds, I will write this as
$\arith(L)^{\mathrm{nor}}_y$'' \hfill\anchor{ATS~II(1/2)}{1867--1873}
\end{quote}

\begin{quote}\small
``Because the absolute values $(|\cdot|_{K_{\psi_v}})$ depend on $y \in \Ynum$,
one has (5.3.7) $\lambda_{y'} \neq \lambda_y$ in general and hence one says that
the \emph{product formula normalization of $y$ moves with $y$} $\in \Ynum$.
Especially, one sees that there is \emph{no uniform choice} of the normalization
coordinate $\lambda_y$ for all $y \in \Ynum$.''
\hfill\anchor{ATS~II(1/2)}{1884--1891}
\end{quote}

Zwei weitere Kontextstücke gehören hierher. Jeder Punkt $y$ trägt eine
\emph{kanonische} Bewertung, gewonnen über die Fargues--Fontaine-Kurve
\cite{FarguesFontaine2018}, und daneben die \emph{$L_v$-standardnormierte} Bewertung
mit $|p|_{K_v} = 1/p^{f_v}$ \anchor{ATS~II(1/2)}{802--840}. Und die Unmöglichkeit,
alles zugleich zu normieren, wird als Aussage über die Familie formuliert:

\begin{quote}\small
``Regardless of the choice of a normalization \ldots\ one cannot normalize
valuations of all untilts of $\bar{L}_v$ \emph{simultaneously}. More precisely the
function $y \mapsto |p|_{K_y}$ is not a constant function.''
\hfill\anchor{ATS~II(1/2)}{802--840, Rem.~2.7.1}
\end{quote}

\subsection{Drei Konsequenzen}

\paragraph{(i) Als stellenweise Bindung gelesen, legt (5.3.3) die Normierung für
jedes feste $y$ eindeutig fest.}
Der Exponent $\lambda_v$ ist durch (5.3.3) als Umrechnungsexponent zwischen der
kanonischen $K_{\psi_v}$-Bewertung und der Artin--Whaples-Standardbewertung
festgelegt. Das ist die klassische Eindeutigkeitsaussage für äquivalente Beträge:
Ist $|\cdot|_1$ nichttrivial und $|\cdot|_2 = |\cdot|_1^{\,a}$ für ein $a > 0$, so
ist $a$ eindeutig, denn es ist als $\log|y|_2/\log|y|_1$ für jedes $y$ mit
$|y|_1 > 1$ berechenbar \cite[Prop.~7.8, S.~108]{Milne2020}. So gelesen gibt es
\emph{keine Restfreiheit}, innerhalb derer irgendetwas \glqq koordiniert\grqq{}
werden könnte. Diese Lesart ist naheliegend --- und das Verdikt dieses Abschnitts
ruht auf ihr ---, aber sie ist nicht die einzige, die die Quellen zulassen;
Abschnitt~\ref{sec:trilemma} formuliert die Alternative.

\paragraph{(ii) Auf Elementen des Zahlkörpers stimmen die normierten Bewertungen
für jedes $y$ mit den Standardbewertungen überein.}
Das ist der Gehalt von (5.3.3)/(5.3.4). Was sich mit $y$ bewegt, ist die
\emph{Koordinate} $\lambda_y$ --- die Lage der Standardskala relativ zur kanonischen
Untilt-Skala (untilt scale) --- und nicht der Wert $|x|$ für $x \in L$.

\paragraph{(iii) Die Nicht-Simultaneität ist eine Aussage über die Familie.}
Bemerkung 2.7.1 verneint eine $y$-uniforme Normierung über $\Ynum$; sie ist kein
Hindernis für die punktweise Normierung eines einzelnen $y$. Die Familienaussage von
der Punktaussage zu unterscheiden, ist genau das, was (i) und das zitierte \glqq no
uniform choice\grqq{} verträglich statt widersprüchlich macht.

\subsection{Entscheidung der Trichotomie, unter der stellenweisen Lesart}
\label{sec:trichotomy}

\begin{itemize}
\item \textbf{N3 ist gegenstandslos} --- wenn (5.3.3) als stellenweise Bindung
  gelesen wird. Abschnitt \S\,5.3 lässt dann keine Restfreiheit: Die in der
  Formulierungsphase vermutete Freiheit --- dass $\prod_v |x|_v = 1$ die
  einzelnen Beträge nicht festlege --- existiert dann nicht, denn
  es wird nicht \emph{irgendeine} produktformelverträgliche Familie gewählt, sondern
  \emph{diejenige}, die durch (5.3.3) an den Artin--Whaples-Standard gebunden ist.
  Abschnitt~\ref{sec:trilemma} zeigt, dass dieser Vorbehalt nicht müßig ist: An
  anderer Stelle der Reihe wird eine Restfreiheit genau der N3-Art
  \emph{ausgeübt}.
\item \textbf{N2 ist unter den Annahmen (F1)--(F2)} der Herleitung in
  Abschnitt~\ref{sec:D1} \textbf{ausgeschlossen}; ohne sie ist es textlich
  unwahrscheinlich, aber offen. Der Ausschluss ist zweigunabhängig --- er benutzt
  allein die Bewegung von
  $\lambda_y$ entlang einer Frobenius-Faser, die jeder Zweig beibehält ---, aber
  er ist nicht annahmefrei: (F1)--(F2) verorten den Standardpunkt aus ATS~IV auf
  einer solchen Faser und liefern die Koordinatentransformation an zwei endlichen
  Stellen.
\item \textbf{N1$'$ gilt unter dieser Lesart und ist dann die Bedeutung der
  Wendung.} Mit N1$'$ meinen wir
  eine geschärfte Form von N1: Beide Kopien werden in ihrer (jeweils eindeutigen)
  normierten Gestalt genommen, wodurch Elemente des Zahlkörpers in beiden Kopien
  dieselben Standardbeträge erhalten. Das ist bei einer stellenweisen Lesart von
  \S\,5.3 das Einzige, was
  \glqq Choose the same global normalization for the number field provided by the
  arithmeticoids $y_0'$ and $y_0$\grqq{} bedeuten kann, und genau das sagt Joshis
  eigene Begründung: Die Arithmeticoide ``provide an isomorphic copy \ldots\ hence
  one has chosen the same global normalization''.
\end{itemize}

Die Formulierung als \glqq Wahl\grqq{} wäre in diesem Licht leicht irreführend: Es
gäbe nichts zu wählen. Sie wäre die Anrufung der stehenden Konvention von ATS~III,
nach der alle Arithmeticoide normiert sind. Frage~\ref{q:main} wäre damit in ihrer
zweiten Alternative beantwortet --- die Wendung stellt die Identität per Definition
her und behauptet keine Isometrie ---, und auf der Ebene der Zahlkörperelemente
existierte ein \glqq fehlendes Vertragslemma\grqq{} der vermuteten Form nicht.

Wir schreiben das mit Bedacht im Konjunktiv. Das Verdikt ist stichhaltig
\emph{gegeben} die stellenweise Lesart von (5.3.3), und diese Lesart ist für den
Wortlaut von \S\,5.3 für sich genommen die naheliegende. Sie ist aber nicht die
Lesart, unter der die Normierung im Rest der Reihe gebraucht wird. Der nächste
Unterabschnitt legt die daraus folgende Alternative dar, und jedes spätere Verdikt
dieser Arbeit ist für die Zweige markiert, unter denen es gilt.

Unabhängig vom Zweig identifiziert die Definition \emph{nichts} auf der
Untilt-Ebene, wo sich die kanonischen Bewertungen benachbarter
Punkte einer Frobenius-Faser nachweislich schon auf $\mathbb{Q}_p$ unterscheiden.
Welche der Größen in Theorem 6.10.1 auf welcher Ebene leben, ist Gegenstand von
Abschnitt~\ref{sec:terms}.

% -----------------------------------------------------------------------------
\subsection{Was die Definition nicht sagt: ein Trilemma für
\texorpdfstring{$\lambda_y$}{lambda\_y}}
\label{sec:trilemma}
% -----------------------------------------------------------------------------

Der vorangehende Unterabschnitt hat eine Antwort auf eine Frage unterstellt, die
die Quellen nicht behandeln: \emph{Welche Regel wählt die Normierungskoordinate
$\lambda_y$ aus, wenn sich das Arithmeticoid bewegt?} Wir haben die in
Abschnitt~\ref{sec:sources} genannten Passagen nach einer solchen Regel abgesucht
und keine gefunden. Gefunden haben wir stattdessen, dass beide Kandidatenantworten
benutzt werden --- an verschiedenen Stellen.

\paragraph{Die stellenweise Lesart ist das, was \S\,5.3 sagt.}
Gleichung (5.3.3) legt $\lambda_v$ als den Exponenten fest, der die kanonische
Bewertung in die Artin--Whaples-Standardbewertung überführt, und (5.3.5) spricht
von \emph{der} Normierungskoordinate. Unter dieser Lesart ist (5.3.4) eine
Umschrift von (5.3.1) und keine zusätzliche Forderung --- der Text sagt selbst,
die Produktformel \glqq reads\grqq{} so.

\paragraph{Eine stellenselektive Freiheit ist das, was die späteren Arbeiten
ausüben.}
ATS~III hält fest, dass die Skalierung an den halbstabilen Stellen eine
Renormierung \emph{an den übrigen Stellen} erzwingt:

\begin{quote}\small
``the non-trivial scaling relationship established by Theorem 4.2.2.1(3) at primes
$w \in V_{\mathrm{odd},ss}$ \emph{forces that suitable (re)normalization of
valuations must be introduced at primes} $w \in V_{L'} - V_{\mathrm{odd},ss}$ so
that the product formula to continue to hold for each normalized arithmeticoid
$\arith(L)^{\mathrm{nor}}_{y'_j}$.'' \hfill\anchor{ATS~III}{1600--1609}
\end{quote}

Derselbe Mechanismus erscheint als Struktureigenschaft in Theorem 4.6.1(4)
\anchor{ATS~III}{1786--1791}, und ATS~IV erklärt die Bewegung der Höhen mit ihm:
Die lokalen Beiträge seien ``\emph{bound together} \ldots\ by means of the
(global) product formula'', sodass der globale Frobenius ``changes global heights
in a subtle and complicated way in order to ensure that the product formula
holds'' \anchor{ATS~IV}{585--589}. Die Anweisung, an anderen Stellen eine
Renormierung einzuführen, setzt voraus, dass dort noch etwas zu wählen ist --- und
genau diese Freiheit bestreitet eine stellenweise Lesart von (5.3.3).

\paragraph{Drei diagnostische Zweige (nicht erschöpfend).}
Da keine Regel angegeben ist, sind drei quellmotivierte Möglichkeiten auseinanderzuhalten:

\begin{description}
\item[(a) Standardgebunden.] $\lambda_y$ \emph{ist} das stellenweise durch (5.3.3)
  bestimmte Tupel. Dann ist $\arith(L)^{\mathrm{nor}}_y$ wohldefiniert, und
  Abschnitt~\ref{sec:trichotomy} gilt wie geschrieben --- aber an den übrigen
  Stellen ist nichts \glqq einzuführen\grqq{}, und die oben zitierte Anweisung hat
  keinen Gegenstand.
\item[(b) Produktformelgebunden.] $\lambda_y$ ist irgendein Tupel, für das die
  Produktformel gilt. Dann ist es nach Lemma~1 unten nur bis auf einen Skalar
  bestimmt, und die Freiheit ist \emph{nicht} stellenselektiv.
\item[(c) Stellenselektiv.] $\lambda_y$ wird durch eine weitere Regel ausgewählt,
  die echte stellenweise Freiheit lässt. Das ist es, was die eben zitierten
  Passagen gebrauchen --- aber die Regel selbst steht nicht im gelesenen Material,
  und mit einer stellenweisen Lesart von (5.3.3) ist sie nicht vereinbar.
\end{description}

Ob (a), (b), (c) oder eine weitere globale Auswahlregel wie die oben genannte
gemeint ist, lässt sich, soweit wir sehen, am Text nicht entscheiden. Es ist die
natürliche erste Frage jeder Korrespondenz über dieses Material, und wir
formulieren sie in Abschnitt~\ref{sec:fneu} als solche.

% -----------------------------------------------------------------------------
\subsection{Was die Produktformel allein liefern kann: ein Starrheitslemma}
\label{sec:rigidity}
% -----------------------------------------------------------------------------

Zweig (b) lässt sich entscheiden. Die Produktformel schränkt das Multiplikatortupel
weit stärker ein, als man erwarten könnte: Sie legt es vollständig fest, bis auf
einen globalen Skalar.

\paragraph{Die vom Lemma verwendeten Umrechnungsmultiplikatoren.}
Sei $\lambda^{\mathrm{std}}_{y,v}>0$ durch
$|x|_{K_{y_v}}^{\lambda^{\mathrm{std}}_{y,v}}=|x|^{\mathrm{std}}_v$ auf $L_v$
bestimmt, und sei $\lambda_{y,v}=\mu_v\lambda^{\mathrm{std}}_{y,v}$.
Dann gilt $|x|_{K_{y_v}}^{\lambda_{y,v}}=(|x|^{\mathrm{std}}_v)^{\mu_v}$.
Das folgende Lemma beschränkt $\mu$, nicht unmittelbar $\lambda_y$.
Zweig~(a) legt $\mu_v=1$ fest; Zweig~(b) beschränkt positive
Umrechnungsexponenten auf $\mu_v=c>0$. Eine entsprechende Aussage über
Haar-Maße wird damit nicht getroffen.

\begin{lemmaR}[Starrheit des Produktformel-Multiplikators]
\label{lem:rigidity}
Sei $L$ ein Zahlkörper, und für jede Stelle $v$ bezeichne
$|\cdot|^{\mathrm{std}}_v$ den Artin--Whaples-Standardbetrag, sodass
$\sum_v \log|x|^{\mathrm{std}}_v = 0$ für alle $x \in L^*$ gilt. Seien
$(\mu_v)_{v \in V_L}$ reelle Zahlen mit
\[
  \sum_{v \in V_L} \mu_v \cdot \log|x|^{\mathrm{std}}_v \;=\; 0
  \qquad \text{für alle } x \in L^*.
\]
Dann sind alle $\mu_v$ gleich.
\end{lemmaR}

\noindent\emph{Beweisskizze.}
Für $L = \mathbb{Q}$ ist die Aussage unmittelbar und braucht nichts über die
Standardbeträge hinaus. Einsetzen von $x = p$ für eine Primzahl $p$ gibt
$\log|p|_p = -\log p$, $\log|p|_\infty = +\log p$ und $\log|p|_q = 0$ für
Primzahlen $q \neq p$; die Voraussetzung lautet damit
$(\mu_\infty - \mu_p)\log p = 0$, also $\mu_p = \mu_\infty$ für jede Primzahl $p$.
$x = -1$ liefert keine weitere Bedingung, und $\mathbb{Q}^*$ wird von $-1$ und den
Primzahlen erzeugt.

Für einen allgemeinen Zahlkörper werden dieselben zwei Schritte geführt, mit zwei
klassischen Zutaten. Erstens verschwinden für eine Einheit $x \in \mathcal{O}_L^*$
alle endlichen Terme, und nach dem Dirichletschen Einheitensatz spannen die
Vektoren $(\log|x|_v)_{v \text{ arch}}$ die Spurnull-Hyperebene der archimedischen
Koordinaten auf; also ist $\mu_v$ über die archimedischen Stellen konstant $= c$.
Zweitens sei für eine endliche Stelle $v$ die Zahl $h$ die Klassenzahl von $L$ und
$\pi$ ein Erzeuger des Hauptideals $\mathfrak{p}_v^{\,h}$; Einsetzen von $x = \pi$
gibt $-\mu_v \cdot h\log N(\mathfrak{p}_v) + c \cdot h\log N(\mathfrak{p}_v) = 0$,
also $\mu_v = c$. \hfill$\square$

\paragraph{Status und Herkunft.}
Das Lemma ist elementar und der Sache nach gewiss nicht neu; wir formulieren es,
weil die von ihm ausgedrückte Einschränkung im Material implizit gebraucht und,
soweit wir feststellen konnten, dort nirgends zitiert wird. Der
$\mathbb{Q}$-Fall oben ist vollständig geführt. Der allgemeine Fall benutzt den
Dirichletschen Einheitensatz und die Endlichkeit der Klassenzahl, beides
klassisch; wir geben für keinen von beiden eine Satznummer an, weil wir in keiner
uns vorliegenden Ausgabe eine verifiziert haben. Ausdrücklich schreiben wir das
Lemma \emph{nicht}~\cite{ArtinWhaples1945} zu: Die dort bewiesene
Produktformel-Charakterisierung betrifft, welche Körper ein solches System von
Beträgen zulassen, und wir konnten den Volltext nicht beschaffen, um zu prüfen, ob
eine Eindeutigkeitsaussage über die Multiplikatoren enthalten ist. Wer dieses
Lemma zitiert, sollte es entweder wie hier beweisen oder eine Quelle unmittelbar
verifizieren.

\paragraph{Konsequenz für Zweig (b), sorgfältig formuliert.}
Nach Lemma~\ref{lem:rigidity} sind die zulässigen Tupel in Zweig (b) genau
$\mu \equiv c$ mit $c > 0$. Daraus folgt zweierlei, und beides zieht in
entgegengesetzte Richtungen.

Erstens ist die Freiheit, die Zweig (b) gewährt, \emph{uniform} und niemals
selektiv. Ein globaler Faktor multipliziert jeden lokalen Beitrag gleichermaßen; er
kann den Beitrag an einer Stelle nicht heben und an einer anderen senken. Wo das
Material eine Renormierung an den nicht-halbstabilen Stellen als \emph{Kompensation}
einer Skalierung an den halbstabilen beschreibt, verlangt diese Beschreibung mehr,
als die Produktformel für sich genommen liefern kann. Das ist eine Aussage darüber,
was eine Nebenbedingung impliziert --- es ist keine Behauptung, Korollar 4.2.2.4 sei
falsch, und keine Behauptung, es gebe keine Regel, die den beschriebenen Effekt
erzielt. Es besagt, dass eine solche Regel von woanders herkommen müsste als von
der Produktformel, und schärft damit, was eine Auflösung des Trilemmas zu liefern
hätte.

Zweitens ist der Skalar $c$ durch die Nebenbedingung nicht festgelegt. Für ein
projektives Objekt wie die Hyperebene $H_y$ ist das harmlos --- $c$ kürzt sich
heraus, da Koeffizientenvektoren genau dann dieselbe Hyperebene definieren, wenn
sie proportional sind ---, für eine in $\lambda$ lineare Größe wie eine Höhe
jedoch nicht: Ihr Wert ändert sich um den Faktor $c$. Zweig (b) lässt die Skala
solcher Größen daher unbestimmt, solange keine weitere Konvention hinzutritt.

\begin{table}[htbp]
\centering
\caption{Welches Verdikt dieser Arbeit unter welchem Zweig des Trilemmas aus
Abschnitt~\ref{sec:trilemma} gilt. \glqq ---\grqq{} markiert ein Verdikt, das der
Zweig in der angegebenen Form nicht trägt. Zweigunabhängigkeit ist nicht dasselbe
wie Annahmefreiheit: Die erste Zeile gilt in jedem Zweig \emph{unter den Annahmen}
(F1)--(F2) aus Abschnitt~\ref{sec:D1}, und die fünfte Zeile benennt den Träger der
mittleren Gleichheit, ohne den Transport zu den beiden in (E2) angeschriebenen
indizierten Hull-Auswertungen herzustellen (Abschnitt~\ref{sec:residues}).}
\label{tab:branches}
\small
\begin{tabular}{@{}P{0.22\textwidth}P{0.15\textwidth}P{0.15\textwidth}P{0.17\textwidth}P{0.10\textwidth}@{}}
\toprule
Verdikt & (a) standard\-gebunden & (b) PF\-gebunden & (c) stellen\-selektiv & \shortstack{Evidenz\\niveau} \\
\midrule
N2 ausgeschlossen (D-1, unter (F1)--(F2)) & gilt & gilt & gilt & \textbf{O} \\[3pt]
N3 gegenstandslos & gilt & teilweise: Freiheit besteht, ist aber nicht selektiv & --- & \textbf{O} \\[3pt]
Skalenidentifikation definitorisch (N1$'$) & gilt & bis auf den Skalar $c$ & --- & \textbf{D} \\[3pt]
E1 ($q_\ell$) invariant & gilt & gilt & gilt (liest keine Skala) & \textbf{D} \\[3pt]
E2, Trägerkandidat ($\varphi$-Stabilität von $\Thet$; indizierter Transport nicht
  nachgezeichnet) & gilt & gilt & gilt & \textbf{S} \\[3pt]
E2-Maßverträglichkeit & offen & offen & offen & \textbf{O} \\[3pt]
Auf der Untilt-Ebene wird nichts identifiziert & gilt & gilt & gilt & \textbf{D} \\
\bottomrule
\end{tabular}
\par\smallskip
{\footnotesize \textbf{D}: definitorisch;\quad \textbf{S}: quellenbehauptet;\quad
\textbf{O}: offen oder bedingt.\par}
\end{table}

% =============================================================================
\section{Die Periodenabbildung bewegt sich wirklich: der konditionale Ausschluss
von N2}
\label{sec:D1}
% =============================================================================

\subsection{Die Beweglichkeit der Hyperebene in den Quellen}

Jedes normierte Arithmeticoid liefert einen Gradhomomorphismus und, über den
Logarithmus der Produktformel, eine Hyperebene:

\begin{quote}\small
``(2) Each normalized arithmeticoid $\arith(L)_y$ provides a degree homomorphism
\ldots\ $(x_v) \mapsto \sum_{v \in V_L} \log(|x_v|_{K_{y_v}}) = \deg_{y_j}((x_v)_{v \in V_L})$.
(3) As a consequence, each normalized $\arith(L)^{\mathrm{nor}}_y$ provides a
hyperplane $H_{y_j} \subset V_L$ given by the logarithm of product formula equation
(2.1.1)'' \hfill\anchor{ATS~III}{1762--1783}
\end{quote}

Die Hyperebene ist nicht fest:

\begin{quote}\small
``the above arithmetic period mapping is a highly non-trivial global arithmetic
period mapping in which an arithmetic avatar $\arith(L)_y$ provides a hyperplane
$H_y$ (as above) \emph{which moves under symmetries (such as the actions of the
absolute Galois group, global Frobenius and $L^*$)}''
\hfill\anchor{ATS~III}{518--522}
\end{quote}

und entsprechend ist ``the global arithmetic period mapping is non-constant''
\anchor{ATS~III}{869--870}. Als Grund wird das Zusammenspiel von lokaler Freiheit und
globaler Nebenbedingung angegeben:

\begin{quote}\small
``the valuations on untilts parameterized by a Fargues--Fontaine curve at some
arbitrary prime \ldots\ cannot be all simultaneously normalized and on the other
hand, local changes of valuations are globally constrained by the validity of the
product formula'' \hfill\anchor{ATS~III}{515--518}
\end{quote}

In der früheren Arbeit wird die Hyperebene \emph{explizit mit den
Normierungskoordinaten als Koeffizienten} angegeben:

\begin{quote}\small
``$H_y = \{ z = (z_v) \in \bigoplus_v \mathbb{R} : \sum_v \lambda_v z_v = 0 \}$''
\hfill\anchor{ATS~II(1/2)}{2199--2206}
\end{quote}

zusammen mit der Aussage, man habe ``a natural and \emph{non-constant} mapping
$y \mapsto H_y$ \ldots\ called the period mapping''
\anchor{ATS~II(1/2)}{2227--2233}, deren Nichtkonstanz als ``immediate from the fact
that \emph{normalizations of arithmeticoids do not carry over} because locally at
each prime $v$, valuations \emph{cannot be simultaneously normalized}'' bewiesen wird
\anchor{ATS~II(1/2)}{2238--2243}. Bemerkung 5.10.2(5) derselben Arbeit hält fest,
dass der Unmöglichkeitseinwand von Scholze--Stix~\cite{ScholzeStix2018} genau dieses
Phänomen betrifft; ATS~III vermerkt ebenfalls, dass Scholze--Stix die Beweglichkeit
von $H_y$ bestreiten \anchor{ATS~III}{522--525}. Wir berichten diese
Meinungsverschiedenheit und beziehen darin keine Position.

\subsection{Das Verhalten entlang der Frobenius-Faser}

Auf der Faser $\{y_n\}_{n \in \mathbb{Z}}$ über dem kanonischen Punkt --- dem Orbit
der Decktransformation, also des Frobenius --- stellen die Quellen fest:

\begin{itemize}
\item die Kummer-Einbettungen $\mathbb{Z}_p(1) \to B$ zu $y_n$ und $y_{n-1}$
  ``differ as $\mathbb{Z}_p$-submodules by a $p$-multiple''
  \anchor{ATS~II}{2119--2120};
\item ``One has $|p|_{K_{y_{n-1}}} = |p|^{1/p}_{K_{y_n}}$''
  \anchor{ATS~II}{2122};
\item ``the valuations of $p$, \emph{and hence of the Tate parameters} of the
  holomorphoids \ldots\ \emph{grow along the fiber} $\{y_n\}$''
  \anchor{ATS~II}{2127--2131};
\item ``each $y_n$ comes with a canonical valuation on its residue field
  $K_{y_n}$ \ldots\ but the valuations provided by $y_{n-1}$ and $y_n$ \emph{are
  not the same on $\mathbb{Q}_p$}'' \anchor{ATS~II}{2184--2188}.
\end{itemize}

und, für die Wirkung an allen Stellen zugleich, dass der globale Frobenius die
lokalen multiplikativen Strukturen von $L$ und $L^{(1)}$ ``through the \emph{$p_v$-th
power mapping for each $v \in V_L$}'' in Beziehung setzt
\anchor{ATS~II(1/2)}{2149--2153}.

\subsection{Die Herleitung}

\begin{derivationD}
Es gelte \textbf{(F1)}, dass der in ATS~IV \S\,4.1 verwendete Standardpunkt $y_0$
in einer Frobenius-Faser liegt, auf die ATS~II Theorem 10.20.1(3) anwendbar ist,
und \textbf{(F2)}, dass ATS~II(1/2) Bemerkung 5.9.2
$\lambda_{\varphi(y),v} = p_v \lambda_{y,v}$ an zwei endlichen Stellen
verschiedener Restcharakteristik induziert. Dann gilt
$H_{\varphi(y)} \neq H_y$: Die Produktformel-Hyperebene wird vom globalen Frobenius
nicht erhalten.
\end{derivationD}

\noindent\emph{Herleitung.}
\begin{enumerate}
\item $H_y$ hat den Koeffizientenvektor $\lambda_y$
  \anchor{ATS~II(1/2)}{2199--2206}.
\item Entlang der kanonischen Faser gilt
  $|p|_{K_{y_{n-1}}} = |p|^{1/p}_{K_{y_n}}$ \anchor{ATS~II}{2122}. Mit der
  Standardbedingung $|p|^{\lambda_n}_{K_{y_n}} = 1/p$ (Fall $L = \mathbb{Q}$,
  $v = p$) ergibt das $\lambda_{n-1} = p \cdot \lambda_n$.
\item An \emph{jeder} Stelle $v$ wirkt $\varphi$ als $p_v$-ter Potenzschritt
  \anchor{ATS~II(1/2)}{2149--2153}, d.\,h.\ $\lambda_v \mapsto p_v \cdot \lambda_v$.
  Der Schritt von der zitierten Aussage --- die von den lokalen
  \emph{multiplikativen Strukturen} spricht --- zu dieser Aussage über die
  \emph{Normierungskoordinaten} ist in Analogie zu Schritt 2 aus dem Zitat gelesen;
  als Rechnung ausgeführt ist er dort nur im Fall $L = \mathbb{Q}$, $v = p$. Das
  ist die Annahme (F2).
\item Zwei Koeffizientenvektoren definieren nur dann dieselbe Hyperebene, wenn sie
  durch einen \emph{einzigen} Skalar $c$ proportional sind. Aber
  $(p_v \lambda_v) = c(\lambda_v)$ erzwänge $p_v = c$ für alle $v$, was unmöglich
  ist, da die Restklassencharakteristiken variieren. Also
  $H_{\varphi(y)} \neq H_y$. \hfill$\square$
\end{enumerate}

\paragraph{Status dieser Herleitung.}
D-1 ist \emph{unsere eigene} Herleitung --- elementare Algebra plus ein
Proportionalitätsargument ---, zusammengesetzt aus oben wörtlich zitierten Aussagen.
Sie ist \emph{kein} Satz der Quelle und wird auch nicht als solcher präsentiert. Sie
steht hier, weil sie Lesart N2 entscheidet, welche die Formulierungsphase nur als
\glqq dokumentiert-unwahrscheinlich\grqq{} einstufen konnte.

\paragraph{Zwei ausdrücklich mitgeführte Vorbehalte.}
\begin{enumerate}
\item[(i)] Schritt 2 rechnet in der kanonischen Faser über dem kanonischen Punkt.
  Dass die in ATS~IV \S\,4.1 fixierten Standardpunkte $y_0$ auf einer solchen Faser
  liegen, ist mit der Aussage verträglich, dass ``there are countably infinite
  possible choices for the standard point $y_0$ \ldots\ Any one of the countably
  many choices for $y_0$ are equally valid choices'' \anchor{ATS~III}{1723--1729},
  wurde hier aber nicht gesondert nachgewiesen.
\item[(ii)] Das Verhalten an den archimedischen Stellen wird in Bemerkung 5.9.2
  durch die Wendung ``for each $v \in V_L$'' behauptet, wurde aber nicht gesondert
  geprüft. Für D-1 wird es nicht gebraucht: Zwei endliche Stellen mit verschiedenen
  Restklassencharakteristiken genügen für Schritt 4 bereits.
\end{enumerate}

Unter diesen beiden Vorbehalten --- den Annahmen (F1)--(F2) der obigen Aussage ---
ist N2 ausgeschlossen und nicht bloß unwahrscheinlich,
und der Fixpunkt-Vorbehalt, den die Formulierungsphase mitführen musste --- die
Möglichkeit, dass $y_0$ zufällig ein $\varphi$-Fixpunkt sein könnte --- entfällt:
Entlang kanonischer Fasern bewegt sich die Normierungskoordinate nachweislich, und
ein fester Standardpunkt müsste außerhalb solcher Fasern liegen, wofür nichts im
Material spricht.

% =============================================================================
\section{Termklassifikation von Theorem 6.10.1: drei Arten von Aussage}
\label{sec:terms}
% =============================================================================

\subsection{Die Kette}

In vereinfachter Notation behauptet Theorem 6.10.1 die Kette
\begin{align}
  -\bigl|\log\bigl(q_\ell^{\varphi(y_0)}\bigr)\bigr| &= -\bigl|\log\bigl(q_\ell^{y_0}\bigr)\bigr|
  \tag{E1}\\[2pt]
  -\bigl|\log\bigl(q_\ell^{y_0}\bigr)\bigr|
  &\le -\tfrac{1}{\ell^*}\bigl|\log \Vol\bigl(\mathrm{hull}(\Thet^{I_M,\varphi(y_0)})\bigr)\bigr|
  \;=\; -\tfrac{1}{\ell^*}\bigl|\log \Vol\bigl(\mathrm{hull}(\Thet^{I,y_0})\bigr)\bigr|
  \;\le\; C_\Theta \bigl|\log\bigl(q_\ell^{y_0}\bigr)\bigr|
  \tag{E2}
\end{align}
\anchor{ATS~IV}{3180--3205}, wobei die erste Ungleichung über das entsprechende
Korollar von ATS~III auf~\cite{Mochizuki2021c} zurückgeführt wird. Unser Gegenstand
ist die mittlere Gleichheit von (E2) und die Gleichheit (E1) --- nicht die
Ungleichungen, die wir nicht untersuchen.

\subsection{E1: \texorpdfstring{$q_\ell$}{q\_l} ist eine klassische, \texorpdfstring{$y$}{y}-freie Größe}

Der $q$-Divisor von ATS~IV \S\,4.4 ist definiert durch

\begin{quote}\small
``let $q_w$ at $C_{M_w}$ be a Tate parameter (\emph{well defined up to
multiplication by a local unit and its order at $w$}). \ldots\
$q_M = \sum_{w \in V^{\mathrm{odd},ss}_M} \ordd_w(q_w) \cdot w$''
\hfill\anchor{ATS~IV}{1831--1846}
\end{quote}

Die vom Theorem verwendete Definitionskette lässt sich kompakt schreiben als
\[
 q_M=\sum_{w\in V_M^{\mathrm{odd},ss}}\ordd_w(q_w)\,w,
 \qquad
 \log(q_M)=\frac{1}{[M:\mathbb{Q}]}\deg(q_M)
 =\frac{1}{[M:\mathbb{Q}]}
   \sum_w \ordd_w(q_w)\log w,
\]
gefolgt unmittelbar vor Theorem~6.10.1 von
$q_\ell^{y_0}=\tfrac{1}{2\ell}\log(q)$, berechnet im Rahmen $y_0$
\anchor{ATS~IV}{1831--1876, 2523--2528, 3180--3189}. Dabei ist $q$ die
Kurzschreibweise aus \S\,5.4 für $q_{C/M}$ und damit für den einschlägigen
\S\,4.4-Divisor $q_M$; entsprechend bezeichnet $\log(q)$ den normierten
arithmetischen Grad $\log(q_M)$ und kein neues Objekt. Der hochgestellte Index
protokolliert damit den Rechenrahmen; der angeschriebene Skalar faktorisiert über den klassischen
Divisor und die Graddaten.

sodass $\ordd_w(q_w)$ die klassische Ordnung der diskreten Bewertung im lokalen
Körper $M_w$ ist; die Klammer erklärt sie zu einer wohldefinierten Invariante der
Definition. Der Beweis von Proposition 4.4.4 ist entsprechend durchweg klassisch:
Uniformisierende $\pi_w, \pi_v$, $q_v = \pi_v^\alpha$,
$\ordd_w(q_w) = e_{w\mid v}\ordd_v(q_v)$, $\log w = f_{w\mid v}\log v$
\anchor{ATS~IV}{1852--1876}. \emph{Kein $y$, kein Untilt und kein Holomorphoid kommt
darin vor.} Der $q$-Divisor ist ein $y$-freies klassisches Objekt des Paares
$(C_M, M)$.

Damit ist E1 \emph{unabhängig vom Trilemma} entschieden, und das ist eigens
festzuhalten, weil es das eine Verdikt hier ist, das keinen Zweig braucht. Die in
$q_\ell$ eingehenden Größen sind ganzzahlige Ordnungen $\ordd_w(q_w)$ und
Restklassenkörper-Mächtigkeiten. Keine von beiden liest einen von einem
Arithmeticoid gelieferten Betrag, und die Reskalierung einer Bewertung mit einem
positiven Exponenten lässt die zugehörige Ordnung unverändert --- eine
Renormierung in irgendeinem der drei Sinne von Abschnitt~\ref{sec:trilemma} kann
$q_\ell$ also nicht bewegen. Welcher Zweig auch gemeint ist:
$q_\ell^{y_0}$ und $q_\ell^{\varphi(y_0)}$ berechnen \emph{dieselbe klassische
Zahl}; der hochgestellte Index hält den umgebenden Rechenrahmen fest, ändert aber
den Wert nicht. Unter Zweig~(a) lässt sich derselbe Schluss über N1$'$ erreichen
--- beide normierten Kopien lesen Zahlkörperelemente in derselben Standardskala
---, doch dieser Weg wird nicht gebraucht, und ihn als Begründung anzuführen ließe
ein unbedingtes Verdikt konditional erscheinen. Die direkte visuelle Rückprüfung
der ATS-IV-Quell-PDF, S.~66, zeigt außerdem Betragsstriche auf \emph{beiden} Seiten
von E1. Die einseitige Fassung im Textextrakt war ein Extraktionsartefakt. Daher
ist keine Zusatzannahme wie $\log(q_\ell)\ge 0$ nötig: E1 ist unter der genannten
Identifikation von $q_\ell$ mit dem $q$-Divisor aus \S\,4.4 definitorisch wahr.

\subsection{Der scheinbare Widerspruch und seine Auflösung}

Auf den ersten Blick kollidiert das mit einer ausdrücklichen Aussage derselben
Arbeit: ``the absolute values of Tate parameters rise under global Frobenius
$\phi$'' \anchor{ATS~IV}{567--569}. Wenn Tate-Parameter steigen, wie kann dann eine
aus ihnen gebildete Größe invariant sein?

Die Auflösung ist, dass zwei \emph{verschiedene Objekte} im Spiel sind. Erstens der
klassische $q$-Divisor von \S\,4.4, der eine Kombination von $\ordd_w$-Werten und
$y$-frei ist --- davon spricht E1. Zweitens das Tate-Idel und die
Untilt-Tate-Parameter, die $z$-abhängig sind; davon sprechen \S\,1.4(2),
Proposition 4.5.12 und die Faseraussagen von ATS~II:

\begin{quote}\small
``Let $\arith(L)^{\mathrm{nor}}_z$ be the \emph{normalized} arithmeticoid \ldots\
the normalized arithmetic degree $\log(TI_{C/\arith(L)_z})$ \emph{is a function of
$z$} \ldots\ one obtains a \emph{non-constant} function `normalized degree of the
Tate divisor' $\Ynum \to \mathbb{R}$'' \hfill\anchor{ATS~IV}{2055--2075}
\end{quote}

Man beachte, dass diese Nichtkonstanz \emph{für das normierte} Arithmeticoid
behauptet wird: Die Normierung beseitigt sie nicht, weil sie überhaupt nicht auf der
Ebene des Zahlkörpers lebt.

Dieselbe Unterscheidung erledigt eine Falle, welche die Formulierungsphase als offen
markiert hatte. Joshi dokumentiert durchaus eine Holomorphoid-Abhängigkeit --- ``an
important point of the theory is that $\mathrm{Tate}(C/M)$ depends on the choice of
a holomorphoid $\hol(C/L)_y$ and one should write $\mathrm{Tate}(C/M)_y$ to indicate
this choice'' \anchor{ATS~IV}{681--684} ---, aber er dokumentiert sie für
$\mathrm{Tate}(C/M)$, und er führt $q$ gerade deshalb ein, weil $\mathrm{Tate}$ eine
Stabilitätseigenschaft fehlt:

\begin{quote}\small
``The divisors $\mathrm{Tate}(C/M)$, $\mathrm{Tate}_2(C/M)$ are not stable under
passage to field extensions \ldots\ so one replaces $\mathrm{Tate}(C/M)_2$ by the
normalized arithmetic degree of another divisor $q = q(C/M)$ (\S\,4.4), which is
similar to the divisor $\mathrm{Tate}$, but which receives contributions only from
primes of $M$ lying over a prime \ldots\ of semi-stable reduction \ldots\ By
Proposition 4.4.4, this new divisor $q$ is stable under base-field extensions.''
\hfill\anchor{ATS~IV}{698--705}
\end{quote}

Die dokumentierte $y$-Abhängigkeit von $\mathrm{Tate}$ auf $q$ zu übertragen, wäre
ein Fehlschluss der Objektvertauschung; schon die Indexmengen unterscheiden sich.
Was die Sache klärt, ist kein Kompensationsmechanismus, sondern die Beobachtung,
dass die $y$-Varianz in einer anderen Formalisierung lebt.

\subsection{E2: Familienstabilität und die offene indizierte Gleichheit}
Die Quell-PDF (ATS~IV v2, S.~66) setzt \emph{Beträge} von Log-Volumina
gleich. Die frühere Wiedergabe ohne diese Striche behauptete einen stärkeren
Sachverhalt. Aus $|a|=|b|$ folgt ohne gemeinsame Vorzeichenbindung nicht
$a=b$; eine solche Bindung ist hier für die genauen indizierten Hüllen
und Volumenkonventionen aus E2 nicht nachgewiesen.

\paragraph{Die Quellenbehauptung: Stabilität der kollationierten Menge.}
ATS~IV beschreibt eine adelische Menge, die Tupel von Tate-Parametern aus
Holomorphoiden kollationiert, und behauptet ihre Stabilität unter Iterationen
des globalen Frobenius \anchor{ATS~IV}{506--515}. Das ist Evidenzniveau~S,
eine Quellenbehauptung. Sie benennt einen möglichen Träger auf Familienebene;
sie identifiziert weder die Teilmengen mit Indizes $I_M,\varphi(y_0)$
und $I,y_0$ noch ihre Umgebungsalgebren oder Maße.
Diese Bindungen wurden hier nicht nachgezeichnet.

\paragraph{Benannter Transport: eine hinreichende stärkere Route, weiterhin offen.}
Sei $T:B_{I_M,\varphi(y_0)}\longrightarrow B_{I,y_0}$ eine vorgeschlagene
Abbildung. Eine mögliche hinreichende Route würde Umgebungsalgebra und
Integralstrukturen, eine Isometrie, die Mengengleichheit
$T(\widetilde\Theta^{I_M,\varphi(y_0)})=\widetilde\Theta^{I,y_0}$,
Hull-Verträglichkeit
$T(\operatorname{hull}(S))=\operatorname{hull}(T(S))$
sowie die Erhaltung der einschlägigen Maße und Gradgewichte nachweisen.
Diese Bedingungen würden gleiche Log-Volumina und damit gleiche Beträge
liefern. Sie sind stärker als die Betragsgleichheit der Quelle und nicht als
notwendig nachgewiesen. Abbildung und sämtliche Verträglichkeiten bleiben
offen. Weder aus dem generischen Hull-Gegenfall noch aus einer speziellen
Algebrazerlegung wird ein Ersatzbeweis für die Maßseite beansprucht.

\paragraph{Die Skalenidentifikation ist nur eine Ebene.}
Unter Zweig~(a) haben Zahlkörperelemente in jeder Kopie die Standardbeträge.
Unter Zweig~(b) beschränkt die Produktformel die Umrechnungsmultiplikatoren
auf einen globalen Skalar, bindet diesen aber nicht zwischen den Rahmen.
Keine dieser Beobachtungen beweist E2 oder transportiert normiertes Haar-Maß.
Unter Zweig~(c) benötigt bereits die Identifikation der Elementskalen eine
eigene Auswahlregel. Die Maßpflichten aus Abschnitt~\ref{sec:measure}
bleiben in allen drei Zweigen offen.

\subsection{Wo die Arbeit sitzt: die Asymmetrie der beiden Schranken}

Eine dritte Textbeobachtung vervollständigt das Bild. Die beiden Schranken der Kette
werden nicht im selben Rahmen berechnet:

\begin{quote}\small
``by the virtue of stability under the global Frobenius morphism, one has the
freedom of the choice of an arithmeticoid in the computing the \emph{lower bound}
\ldots\ [it] is \emph{independent of the choice} \ldots\ the lower bound will be
computed using $y_0' = \varphi(y_0)$ while the upper bound will be computed using
$y_0$ (\emph{the upper bound is not independent of this choice}).''
\hfill\anchor{ATS~IV}{733--742}
\end{quote}

Die Quelle hält außerdem fest, dass eine frühere Fassung diese
Frobenius-Verschiebung ausgelassen hatte \anchor{ATS~IV}{755}.
Das benennt eine mögliche Rolle der Rahmenwahl, entscheidet die Gleichheit
aber nicht. Zweig~(a) macht die Identifikation der Elementskalen definitorisch,
während Verschiebung, indizierte Mengen, Hüllen und Maße weiterhin
verträgliche Bindungen benötigen. Zweig~(c) lässt zusätzlich die
Auswahlregelfrage offen. Diese Einordnung verifiziert weder eine der
Ungleichungen noch zertifiziert sie, dass Familienstabilität den fehlenden
indizierten Schritt liefert.

\subsection{Klassifikationstabelle}

Tabelle~\ref{tab:terms} hält die Klassifikation fest. Zwei Zeilen sind bewusst nicht
in eine der Invarianzklassen gezwungen: $\Thet$ als Menge ist keine skalarwertige
Größe, die als invariant oder nicht invariant klassifiziert werden könnte, und
$\log \Vol(\mathrm{hull}(\Thet))$ ist das Kompositum, dessen Analyse Gegenstand von
Abschnitt~\ref{sec:measure} ist.

\begin{table}[htbp]
\centering
\caption{Termklassifikation von Theorem 6.10.1 mit zugeordneten Evidenz-Niveaus. Klasse~T: definitorisch/trivial
invariant (faktorisiert über $\varphi$-fixierte Daten). Klasse~N: nachweislich
nicht invariant (liest Untilt-Daten). Evidenz-Niveaus: \textbf{D} (definitorisch), \textbf{S} (quellenbehauptet), \textbf{T} (Transport verifiziert), \textbf{O} (offen oder bedingt). Die Klassifikation ist unabhängig vom Trilemma aus
Abschnitt~\ref{sec:trilemma}; in der letzten Zeile bleiben indizierter Transport
und frameübergreifende Maßverträglichkeit in jedem Zweig offen
(Tabelle~\ref{tab:branches}). Klasse~T und Evidenzniveau~T sind verschiedene Kürzel.}
\label{tab:terms}
\small
\setlength{\tabcolsep}{3.5pt}
\begin{tabular}{@{}P{0.23\textwidth}P{0.21\textwidth}P{0.07\textwidth}P{0.11\textwidth}P{0.25\textwidth}@{}}
\toprule
Term & Faktorisiert über & Klasse & \shortstack{Evidenz\\niveau} & Verhalten \\
\midrule
$\ordd_w(q_w)$ (\S\,4.4)
  & Struktur der diskreten Bewertung von $M_w$ ($\varphi$-fixiert)
  & T & \textbf{D} & invariant \\[3pt]
$\log w$, $[M:\mathbb{Q}]$, $e_{w\mid v}$, $f_{w\mid v}$
  & Restklassen- und Graddaten ($\varphi$-fixiert)
  & T & \textbf{D} & invariant \\[3pt]
$q_\ell = \tfrac{1}{2\ell}\log(q)$
  & nur das Obige
  & \textbf{T} & \textbf{D} & \textbf{E1 definitorisch}; Quell-PDF mit Beträgen auf beiden Seiten \\[3pt]
Grad des Tate-Idels $\log(TI_{C/\arith(L)_z})$
  & Untilt-Bewertungen
  & \textbf{N} & \textbf{S} & nicht konstant (Prop.\ 4.5.12) \\[3pt]
Kummer-Einbettung $\mathbb{Z}_p(1) \to B$
  & Wahl des Untilts $y_n$
  & N & \textbf{S} & unterscheidet sich um ein $p$-Vielfaches (Thm.\ 10.20.1(2)) \\[3pt]
$\Thet_{\mathrm{Mochizuki}}$ als Menge
  & die ganze Familie (kollationiert)
  & --- & \textbf{S} & $\varphi$-stabil konstruktionsgemäß \\[3pt]
$\log \Vol(\mathrm{hull}(\Thet))$
  & Menge $+$ Maßkonvention
  & E2 & \textbf{O} & Betrags-E2: indizierter Transport und Maßverträglichkeit in jedem Zweig offen \\
\bottomrule
\end{tabular}
\end{table}

Kriterium~K hat illustrative Beispiele im klassischen $q$-Divisor und
den Untilt-abhängigen Termen. Diese Beispiele machen es weder zu einem
allgemeinen Invarianzsatz noch den algebraischen Hull oder das gewichtete
frameübergreifende Log-Volumen zu transportverifizierten Größen.

% =============================================================================
\section{Maßkonventionen und offene frameübergreifende Pflichten}
\label{sec:measure}

\paragraph{Absolutbetrag, Haar-Maß und Multiplikationsmodul.}
Für jede Stelle und jeden Rahmen unterscheiden wir den Absolutbetrag
$|\cdot|_{y,v}$, das additive Haar-Maß $m_{y,v}$ mit der Normierung
$m_{y,v}(\mathcal O_{L'_w})=1$ und seinen Multiplikationsmodul
\[
 \Delta_{y,v}(\alpha)=\frac{m_{y,v}(\alpha A)}{m_{y,v}(A)}
\]
für eine messbare Menge $A$ positiven endlichen Maßes und $\alpha\ne0$.
Die Gleichheit $\Delta_{y,v}(\alpha)=|\alpha|_{y,v}$ verlangt eine
verträgliche Betragskonvention. Auf $\mathbb Q_p$ erzwingt die Normierung
\[
 m(\mathbb Z_p)=1,\quad m(p\mathbb Z_p)=p^{-1},\quad |p|_p^c=p^{-c}.
\]
Die zweite Identität folgt aus der Zerlegung von $\mathbb Z_p$ in seine
$p$ Nebenklassen modulo $p\mathbb Z_p$. Ein geänderter Bewertungsexponent
ändert dieses normierte additive Haar-Maß nicht.
Der Skalar aus Lemma~\ref{lem:rigidity} liefert daher keinen
Maßtransportbeweis.

\paragraph{Lokale Volumenkonventionen und Gradgewichte.}
ATS~III definiert normiertes Haar-Maß auf den Zahlkörpervervollständigungen
durch $\Vol_{L'_w}(\mathcal O_{L'_w})=1$
\anchor{ATS~III}{6562--6572} und verknüpft geeignete Volumina von
Nebenklassen gebrochener Ideale mit lokalen Absolutbeträgen
\anchor{ATS~III}{6599--6620,6663--6676}.
Die Quelle legt Kehrwerte lokaler Grade als Gewichte fest
\anchor{ATS~III}{6678--6688}.
Das sind Aussagen über festgelegte Körper, Integralstrukturen und
Maßkonventionen. Ihre numerischen Bestandteile sind klassisch; sie
identifizieren weder Umgebungs-Tensoralgebren noch genaue Theta-Mengen,
Hüllen, gewichtete Komponenten oder Maße zwischen den E2-Rahmen.
Standard-Elementskalen unter Zweig~(a) oder ein gemeinsamer
Bewertungsexponent unter Zweig~(b) genügen nicht.
Zweig~(c) lässt zusätzlich die Wahl der Elementskala offen.

\paragraph{Algebraischer Hull muss kein konvexer Abschluss sein.}
ATS~III v4, \S\S\,9.10.7--9.10.8 (S.~125--127), definiert den holomorphen
Hull über Ganzheitsringe der Körperkomponenten der Tensoralgebra.
Die frühere generische Gleichsetzung mit dem konvexen Abschluss war falsch.
Seien $E/\mathbb Q_p$ unverzweigt quadratisch, $E_1=E$,
$E_2=\mathbb Q_p$, die Algebraabbildung die Identität, und $U=\mathbb Z_p\cdot1$.
Die Menge ist kompakt und bereits $\mathbb Z_p$-konvex. Der kleinste
Hull der vorgeschriebenen Form $\alpha\mathcal O_E$, der $1$ enthält, ist
$\mathcal O_E$: Aus $1\in\alpha\mathcal O_E$ folgt
$\alpha\mathcal O_E\supseteq\mathcal O_E$, mit Gleichheit für
$\alpha=1$. Somit gilt
\[
 \operatorname{conv}(U)=\mathbb Z_p\cdot1
 \subsetneq\mathcal O_E=\operatorname{hull}(U).
\]
Die echte Inklusion folgt aus den $\mathbb Z_p$-Rängen eins und zwei.
Die Konvexität einer Hülle, die $U$ enthält, liefert allgemein nur
$\operatorname{conv}(U)\subseteq\operatorname{hull}(U)$.
Die Zulässigkeit dieses $U$ als spezielle Mochizuki-Tensorregion ist
nicht nachgewiesen. Der Gegenfall widerlegt unsere generische Identifikation,
nicht Proposition~9.10.8.1(3), Theorem~9.11.1 oder $abc$.
Ein vorgeschlagener Ersatz über die Zerlegung endlicher étaler Algebren
benötigt weiterhin konkrete Algebraabbildungen, Integralstrukturen,
Mengen, Gewichte und indizierten Transport; er ist kein verifiziertes
positives Maßargument.

\paragraph{Revidiertes Fazit.}
Die frühere Behauptung, die Maßseite faktorisiere restlos über Zahlkörperdaten
oder sei unter Zweig~(a) festgelegt und unter Zweig~(b) bis auf einen Skalar,
wird zurückgenommen. Frameübergreifende Maßverträglichkeit bleibt in jedem
Zweig offen. Identifizierte Elementskalen schließen indiziertes E2 nicht.

\begin{table}[htbp]
\centering
\caption{Komponenten nach der Korrektur. D: definitorisch; S: quellenbehauptet;
O: offen oder bedingt.}
\label{tab:resolution}
\small
\begin{tabular}{@{}P{0.23\textwidth}P{0.27\textwidth}P{0.08\textwidth}P{0.29\textwidth}@{}}
\toprule
Komponente & Grundlage & Niveau & Status \\
\midrule
Elementskalen & Stellenweise Standardlesart & O
 & Definitorisch unter (a); globaler Skalar unter (b); offene Auswahl unter (c) \\[3pt]
E1 & Klassischer $q$-Divisor & D
 & Unter der angegebenen Identifikation; beidseitige Beträge \\[3pt]
Kollationierte Familie & Frobenius-Stabilitätsbehauptung & S
 & Trägerkandidat; indiziertes E2 nicht geliefert \\[3pt]
Indiziertes E2 & Beträge von Log-Volumina & O
 & In jedem Zweig offen; betragslose Gleichheit nicht gefolgert \\[3pt]
Maß und Hull & Haar-Moduln, Integralstrukturen, algebraischer Hull, Gradgewichte & O
 & Verträglichkeit in jedem Zweig offen; generische Hull-Gleichheit zurückgenommen \\[3pt]
Auswahlregel & Koordinate $\lambda_y$ & O
 & Offene Korrespondenzfrage \\[3pt]
Lokaler Vergleich & ATS~IV Bemerkung~6.10.2 & S
 & Ausdrücklich ausgeschlossen \\
\bottomrule
\end{tabular}
\end{table}

% =============================================================================
\section{Woher das Trilemma kommt: die Höhen-Passage}
\label{sec:fneu}
% =============================================================================

Das Trilemma aus Abschnitt~\ref{sec:trilemma} gehörte nicht zu der Frage, die diese
Arbeit beantworten wollte. Es ging aus einer beim Lesen angetroffenen Passage
hervor, und dieser Abschnitt hält fest, wie --- denn das Leitprinzip dieses
Programms lautet, zu dokumentieren, was gesehen wurde, und nicht, eine Position zu
verteidigen.

\begin{findingF}
Die Höhenpassagen aus ATS~II(1/2) \S\S\,8.2--8.3 sind Indizien einer
möglichen Spannung unter der stellenweisen Lesart. Behauptete Nichtkonstanz,
ihr Beweis mit einer lokalen Bewertung und ihr Definitionsbereich
$\mathbb P^n(\overline L)$ sind getrennt zu prüfen.
Eine Auswahlregel für $\lambda_y$ kann die Lesart klären; Domänen- und
Beweisfragen sind durch diese Arbeit aber nicht entschieden.
F-NEU-1 bleibt offen.
\end{findingF}

Die Definition verwendet ausdrücklich die \emph{normierten} Bewertungen: ``Choose the
standard normalization \ldots\ work with $\arith(L)^{\mathrm{nor}}_y$ \ldots\ work
with the normalization coordinate $\lambda_y$'', was
$h(P) = \sum_v \max_j \log\bigl(|x_j|^{\lambda_v}_{K_{y_v}}\bigr)$ ergibt
\anchor{ATS~II(1/2)}{2947--2960}; und Lemma 8.2.3 leitet die Unabhängigkeit der Höhe
von der Wahl homogener Koordinaten gerade daraus ab, ``that we work with a
normalized arithmeticoid'' \anchor{ATS~II(1/2)}{2982--2986}. Unter einer
stellenweisen Lesart von (5.3.3) stimmen diese normierten Werte auf Elementen von
$L_v$ mit den Standardwerten überein --- für \emph{jedes} $y$ ---, sodass die Höhe
für $P \in \mathbb{P}^n(L)$ die klassische, $y$-unabhängige Höhe im
Sinne von~\cite{BombieriGubler2006} wäre. ATS~IV sagt eben dies über die
Definition: ``except for explicating the dependence on the arithmeticoid, the
definition of height \ldots\ is the same as the one given in [Bombieri and Gubler,
2006, \S\,1.5.1]'' \anchor{ATS~IV}{548--550}.

Proposition 8.3.1 behauptet jedoch, dass
``$h_{\arith(L)}$ depends \emph{non-trivially} on
$\arith(L)$'' \anchor{ATS~II(1/2)}{2990--3004}, und ihr Beweis beruft sich darauf,
dass ``$y \mapsto |p|_{K_y}$ is a non-constant function'' --- was die
\emph{kanonische}, unnormierte Bewertungsfunktion aus Bemerkung 2.7.1 ist. Bemerkung
8.3.2 ergänzt, dass ``the
normalization of $\arith(L)^{\mathrm{nor}}_{y_1}$ \emph{cannot be carried over} to
$\arith(L)^{\mathrm{nor}}_{y_2}$'' \anchor{ATS~II(1/2)}{3008--3015} --- womit der
eine einfache Ausweg verschlossen ist, nämlich $\lambda_{y_1}$ festzuhalten, während
sich $y$ bewegt.

Drei mögliche Lesarten verlangen getrennte Prüfungen. Die Definition
in ATS~II(1/2) \S\,8.2 und Lemma~8.2.3 verwendet eine normierte
Höhe; ihre Beziehung zu kanonischen lokalen Bewertungen darf nicht durch
Ersetzung eines Objekts durch ein anderes entschieden werden.
Insbesondere behauptet Proposition~8.3.1 Nichtkonstanz dieser normierten
Höhe auf $\mathbb P^n(\overline L)$, obwohl ihr Beweis auch einen
unnormierten lokalen Absolutbetrag verwendet. Behauptung, Beweisschritt und
Definitionsbereich bleiben getrennte Prüfpflichten.

Eine Beschränkung der $y$-Abhängigkeit auf $\overline L\setminus L$
kann nicht allein aus ATS~IV:545 verworfen werden: Die Original-PDF v2,
S.~13, schreibt $x\in\overline L^*$, nicht $x\in L^*$.
Die Anwendung auf ganzzahlige $abc$-Tripel ist eine getrennte Passage
\anchor{ATS~IV}{608--616}. Unabhängig davon verwendet ATS~II(1/2)
v12, Theorem~8.7.1 (S.~54--55), tatsächlich $x\in L^*$ und
behauptet allgemeine Striktheit \anchor{ATS~II(1/2)}{3210--3216}.
Das ist ein eigenes Textindiz und kein Ersatz für das falsche ATS~IV-Zitat.
Die Familienaussage, eine Normierung könne nicht übertragen werden,
erzeugt allein keine Striktheit an einem festen Punkt. F-NEU-1 bleibt
daher eine offene Korrespondenzfrage; keine dieser Beobachtungen liefert
einen globalen ATS-Fehlernachweis.

Was die Passage dagegen liefert, ist der Mechanismus --- und er ist es, der das
Trilemma erzeugt. In \S\,8.5 wird die Bewegung der Höhe aus der
Produktformel-Nebenbedingung abgeleitet:

\begin{quote}\small
``since one is working with normalized arithmeticoids $\arith(L)^{\mathrm{nor}}$,
the normalization constraint i.e.\ the product formula (5.3.4) implies that
\emph{changes in $p$-adic valuations forced by the action of $n$ must be
compensated by changes at other primes} so that the normalization constraint
\ldots\ continues to hold.'' \hfill\anchor{ATS~II(1/2)}{3162--3193}
\end{quote}

Kompensation an anderen Stellen setzt Freiheit an diesen Stellen voraus. Die
operative Bedeutung von $\arith(L)^{\mathrm{nor}}$ ist in dieser Passage also nicht
\glqq stellenweise an den Artin--Whaples-Standard gebunden\grqq{}, sondern
\glqq global durch die Produktformel eingeschränkt\grqq{} --- und dann lautet die
Frage, welche Regel die Koordinate aus der zulässigen Menge auswählt. Das ist Zweig
(b) oder (c) aus Abschnitt~\ref{sec:trilemma}, und Lemma~\ref{lem:rigidity} besagt,
dass (b) nicht leistet, was die Passage braucht, denn die Freiheit, welche die
Produktformel lässt, ist ein einziger globaler Skalar und kann nicht stellenselektiv
sein.

\paragraph{Die daraus folgende Frage und ihr Status.}
Die genaue Frage, die wir dem Autor vorlegen würden, lautet demnach: \emph{Welche
Regel wählt die Normierungskoordinate $\lambda_y$ für ein bewegtes Arithmeticoid
aus?} Konkret --- ist das normierte Arithmeticoid dasjenige, dessen induzierte
Beträge auf $L_v$ die Artin--Whaples-Standardbeträge sind, wie (5.3.3) nahelegt,
oder ist $(\cdot)^{\mathrm{nor}}$ eine Nebenbedingung an eine Wahl, für die die
Produktformel allein nicht ausreicht, um $\lambda_y$ festzulegen?

Wir erheben keine Fehlerbehauptung, weder gegen Proposition 8.3.1 noch gegen
sonst etwas. Die auflösende Passage kann sehr wohl in Material stehen, das wir nicht
gelesen haben, oder aus einer Lesart der Definition folgen, die der Autor sofort
liefern würde. Die beiden obigen Textbeobachtungen --- das stellenweise Auftreten
von (5.3.3) und der stellenselektive Gebrauch der Normierung in \S\,8.5, in ATS~III
Korollar 4.2.2.4 und in \S\,1.5 von ATS~IV --- sind als Beobachtungen berichtet. Was
sich gegenüber der ersten Fassung dieser Einordnung geändert hat, ist allein dies:
Die Spannung lässt sich nicht als randständig beiseitelegen, weil die Verdikte der
Abschnitte~\ref{sec:trichotomy}--\ref{sec:measure} davon abhängen, wie sie
aufgelöst wird. Diese Abhängigkeit hält Tabelle~\ref{tab:branches} fest. Die diesem
Abschnitt zugrunde liegende Lektüre der Quellen ist in der Begleitprüfung dieses
Programms ausführlicher durchgeführt worden (Teile V und VI, in Vorbereitung).

% =============================================================================
\section{Einordnung und Konsequenzen}
\label{sec:assessment}
% =============================================================================

\subsection{Für die Ledger-Methode}

Der ursprüngliche Ledger lokalisierte eine Frage an der Wendung, die
die Normierungsidentifikation als Tautologie bezeichnet. Unter der
stellenweisen Lesart identifiziert diese Wendung tatsächlich die
Elementskalen des Zahlkörpers definitorisch. Dieser begrenzte Schluss
erledigt E2 nicht: Die Stabilitätsbehauptung der Quelle betrifft eine
kollationierte Familie, während der Satz zwei indizierte algebraische
Hüllen mit gewichteten Volumenkonventionen auswertet.

Die korrigierte Prüfung trennt diese Pflichten. Die Familienbehauptung wird
mit Evidenzniveau~S berichtet; indizierte Gleichheit, Vorzeichenbindung
und frameübergreifende Maßverträglichkeit bleiben auf Niveau~O.
Die generische Hull-Identifikation des früheren positiven Maßarguments
wird zurückgenommen. Unter Zweig~(c) ist die Auswahl der
Normierungskoordinate eine weitere offene Pflicht.
Ein formal abgenommener Ledger oder eine Quellenstellenprüfung beweist
diese Pflichten nicht und stärkt keine ungeprüfte Schranke.

Der nutzbare Ertrag der Methode sind hier die Korrektur konkreter
Quellenlesungen und präzisere offene Pflichten. Er ist kein vollständiges
Brückenzertifikat und keine Abnahme des gesamten ATS-Arguments.

\begin{table}[htbp]
\centering
\caption{Ausgewählte Arbeiten der Gap-Serie (nicht vollständiger Quervergleich): Zielaussagen, primäre Diagnoseinstrumente und Kernergebnisse zu Normierung und logischem Transport.}
\label{tab:gap_series_crosswalk}
\small
\setlength{\tabcolsep}{3.5pt}
\begin{tabular}{@{}P{0.14\textwidth}P{0.20\textwidth}P{0.22\textwidth}P{0.25\textwidth}P{0.11\textwidth}@{}}
\toprule
Teil / Manuskript & Zielaussage & Primäres Instrument & Kernergebnis & Status \\
\midrule
\textbf{Teil~I}\newline (Methode)
  & Inter-universeller Transport
  & Bridge-Contract-Kriterium (C1--C7)
  & Methodisches Fundament; formuliert Verifikationspflichten
  & Veröffentl. \\[3pt]
\textbf{Teil~IV}\newline (Diese Arbeit)
  & Globale Normierung in ATS~IV Thm~6.10.1
  & Invarianz-Klassifikation \& Trilemma-Taxonomie
  & Identifikation ist definitorisch unter (a); Transport der $\Thet$-Hülle unnachgezeichnet
  & Reparatur\newline v1.2 \\[3pt]
\textbf{Teil~IV-B}\newline (BrXi-Anw.)
  & Selbstanwendung von $\mathrm{Br}_{\Xi}$
  & Sub-Bridge Ledger-Audit
  & Geltungsbereich-Einschränkung; verhindert zyklische Domänenvererbung
  & Entwurf v0.1 \\[3pt]
\textbf{Teil~V}\newline (Substanz)
  & Normierungs\-gewichte in ATS
  & Allstellen-Starrheits\-klassifikation
  & Schließt nicht-triviale stellenselektive Skalierung unter Artin--Whaples aus
  & Entwurf v0.1 \\[3pt]
\textbf{Teil~VI}\newline (Reparatur)
  & Reparatur\-programm der ATS-Kette
  & Effektiver Gewichtstest \& Falltaxonomie
  & Beweist Allstellen-Starrheitssatz für Gewichte; testet Höhen-Subkette
  & Entwurf v0.2 \\
\bottomrule
\end{tabular}
\end{table}

\subsection{Für rein globale Mechanismen: ein vorsichtiger Ausblick}
\label{sec:b2outlook}

Das Strukturmuster von Theorem 6.10.1 --- \glqq kein Vergleich Stelle für Stelle,
global festgemacht\grqq{} --- ist dieser Konstruktion nicht eigentümlich. Drei
Träger dieses Musters treten im Material auf, zwei davon aus Joshis eigenem Text.

Der Prototyp ist die Produktformel selbst~\cite{ArtinWhaples1945}. Die Identität
$\prod_v |x|_{K_v} = 1$ sagt über die einzelnen $|x|_v$ überhaupt nichts und ist
gleichwohl die tragende globale Nebenbedingung der gesamten Konstruktion. Der zweite
Träger ist Joshis eigenes zweites Beispiel innerhalb derselben Beweispassage:
``$A' = C \cdot A$. \emph{This is a global equation, which only appears after one
sums over all primes}'' \anchor{ATS~IV}{834--836}. Dass zwei solche Aussagen in einem
Argument nebeneinanderstehen, legt ein wiederkehrendes Muster nahe statt einer
Eigenheit der mittleren Gleichheit. Der dritte ist das in dieser Arbeit untersuchte
Muster, jetzt als ausgearbeitetes Exemplar verfügbar: Die Hyperebenen-Konstruktion
plus Bemerkung 6.10.2 geben einen Fall, in dem die globale Festmachung und die lokale
Unvergleichbarkeit beide ausdrücklich und beide belegt sind.

Eine Arbeitshypothese des vorliegenden Programms --- in seinem eigenen
Kandidatenregister festgehalten und hier mit der gebotenen Vorsicht formuliert ---
ist, dass solche Aussagen auf der kollationierten Familie leben statt an der
einzelnen Stelle, und dass ihnen das eine strukturelle Silhouette gemeinsam mit
Aussagen gibt, die erst nach Summation über alle Stellen verfügbar sind, die
Weil-Positivität darunter. Wir formulieren dies als Hypothese über die \emph{Gestalt
von Argumenten}. Es ist ausdrücklich \emph{keine} mathematische Verbindung und
\emph{keine} Aussage über die Riemannsche Vermutung oder über irgendeine ihrer
Folgerungen. Das zugehörige Arbeitspaket trägt eine ausdrückliche Abbruchklausel:
Sollte sich zeigen, dass die gesammelten Träger nichts über die Beobachtung selbst
hinaus teilen, ist die Analogie als negatives Ergebnis festzuhalten und der
Registereintrag herabzustufen --- nicht stillschweigend weiterzuführen. Eine Frage
in Gegenrichtung im Umfeld von~\cite{GranvilleStark2000} ist im Konzeptdokument des
Projekts vermerkt und wird hier ausdrücklich nicht aufgegriffen.

\subsection{Quellenkorrekturen und offene Pflichten}
\label{sec:residues}

\paragraph{Drei korrigierte Lesungen.}
Die Quell-PDF nennt in ATS~IV v2, S.~13, die Domäne $\overline L^*$.
ATS~II(1/2) Theorem~8.7.1 verwendet getrennt davon $L^*$;
diese Anker sind nicht austauschbar (Abschnitt~\ref{sec:fneu}).
ATS~IV v2, S.~66, setzt in E2 Betragsstriche um beide Log-Volumen-Terme.
Ein gemeinsames Vorzeichen der genauen indizierten Auswertungen ist hier
nicht bewiesen; die frühere betragslose Gleichheit wird deshalb zurückgenommen.
ATS~III v4, S.~125--127, verwendet einen algebraischen Hull;
dessen generische Gleichsetzung mit dem konvexen Abschluss widerlegt
das explizite Beispiel in Abschnitt~\ref{sec:measure}.
Das entscheidet die speziellen Tensorregionen-Aussagen nicht.
Die frühere E1-Vorzeichendiskrepanz bleibt getrennt davon als
Extraktionsartefakt geklärt: Die Quelle setzt in E1 beidseitige Betragsstriche.

\paragraph{Indizierte Gleichheit und Maße (offen).}
Familienstabilität ist quellenbehauptet. Eine Abbildung zwischen den
genauen indizierten Umgebungsalgebren mit ihren Integralstrukturen,
Theta-Mengenbildern, Hüllen, Maßen und Gewichten wurde nicht nachgezeichnet.
Der benannte Transport ist nur eine hinreichende stärkere Route zu E2;
er ist keine aus der Quelle gefolgerte notwendige Bedingung.
Weder eine Normpotenz noch eine generische Hull-Identität liefert das
fehlende Maßargument. Diese Pflichten bleiben in jedem Zweig offen.

\paragraph{Auswahl- und Ordnungsfragen.}
Die Auswahlregel für $\lambda_y$ bleibt offen.
Lemma~\ref{lem:rigidity} beschränkt nur Multiplikatoren von
Standardbewertungen mit Produktformel; es liefert weder eine
stellenselektive Regel noch einen Haar-Maßtransport.
D-1 bleibt an (F1)--(F2) gebunden. F-NEU-1 und die starke
Ordnungsfrage L-strong bleiben offen.
Getrennte Folgeprüfungen sind erforderlich, bevor positive E2- oder
Maßaussagen aus den Begleitteilen~V und~VI als verifizierte Ergebnisse
importiert werden können.

\subsection{Grenzen der Arbeit}
\label{sec:limitations}

Dieser Unterabschnitt benennt, was die Arbeit nicht leistet, in den Worten, in denen
die zugrunde liegenden Arbeitsnotizen es benannt haben.

\begin{enumerate}
\item \textbf{Kein Beweis wurde nachvollzogen.} Weder Theorem 6.10.1 von ATS~IV noch
  das Korollar von ATS~III, auf dem seine erste Ungleichung beruht, noch Mochizukis
  Korollar~3.12~\cite{Mochizuki2021c} wurde Zeile für Zeile verifiziert. Keine
  Abschätzung wurde nachgerechnet und keine Schranke geprüft. Klassifiziert wurden
  die \emph{Terme und die Struktur der Begründung}, nicht die Korrektheit der
  Ungleichungen.
\item \textbf{D-1 ist unsere eigene Herleitung und ist konditional}, kein Satz der
  Quelle: elementare
  Algebra plus ein Proportionalitätsargument, zusammengesetzt aus wörtlich zitierten
  Aussagen. Sie gilt unter den mit ihr genannten Annahmen (F1)--(F2) --- der Annahme
  der kanonischen Faser für den Standardpunkt und der Übertragung der
  Koordinatentransformation auf zwei endliche Stellen, die aus dem Zitat gelesen und
  dort nicht gerechnet ist. Der archimedische Fall ist ebenfalls ungeprüft und für
  D-1 nicht nötig.
\item \textbf{Verdikte über Elementskalen sind zweigkonditional.}
  Definitorische Identifikation und Gegenstandslosigkeit von N3 gelten unter
  Zweig~(a); Zweig~(b) lässt einen globalen Umrechnungsskalar und Zweig~(c)
  seine Auswahlregel offen. Frameübergreifende Maßfestlegung ist in keinem
  Zweig bewiesen; ihre frühere positive Klassifikation wird zurückgenommen.
  D-1 benötigt (F1)--(F2); E1 behält seine angegebene Identifikation;
  Familienstabilität ist quellenbehauptet und kein indizierter Hull- oder
  Volumenbeweis.
\item \textbf{Das Trilemma ist eine Definitionslücke, keine Fehlerbehauptung.} Wir
  haben in den gelesenen Passagen keine Auswahlregel für $\lambda_y$ gefunden; wir
  behaupten nicht, dass es keine gibt, und wir behaupten nichts über die Korrektheit
  von ATS~III Korollar 4.2.2.4 oder von irgendetwas, das darauf ruht.
\item \textbf{Lemma~\ref{lem:rigidity} steht unter einer bewussten
  Referenzbeschränkung.} Der $\mathbb{Q}$-Fall ist hier vollständig bewiesen; der
  allgemeine Fall zieht den Dirichletschen Einheitensatz und die Endlichkeit der
  Klassenzahl als klassische Resultate ohne Satznummer heran, weil wir keine
  verifiziert haben. Insbesondere schreiben wir das Lemma
  nicht~\cite{ArtinWhaples1945} zu, dessen Volltext wir nicht beschaffen konnten.
\item \textbf{F-NEU-1 bleibt offen.} Die korrigierte ATS~IV-Domäne
  $\overline L^*$ und die getrennte ATS~II(1/2)-Aussage auf $L^*$ entscheiden
  weder das gesamte Höhenargument noch einen Fehler in ATS.
\item \textbf{Die Klassifikation von E1 als definitorisch} gilt modulo der Lesart,
  dass $q_\ell$ in Theorem 6.10.1 der
  $q$-Divisor aus \S\,4.4 ist. Die unmittelbar vor dem Theorem gegebene Definition
  stützt diese Lesart, aber die Identifikation ist unsere. Die Quell-PDF selbst
  zeigt auf beiden Seiten von E1 Betragsstriche.
\item \textbf{Die Maßanalyse klassifiziert Konventionen, keine Abschätzungen.} Die
  Korrektheit der Volumenabschätzung selbst --- die Produktschranke, die
  $\ell^*$-Faktoren, die Übersetzung in Mochizukis Formulierung --- wurde nicht
  untersucht.
\item \textbf{Die starke Form der Frage ist unentwickelt.} Die Frage der
  Ordnungsverträglichkeit (L-strong) aus Abschnitt~\ref{sec:equivariance} benennt eine
  Struktur, aber keine Kandidatenklasse. Dies ist der größte inhaltliche Rückstand.
\item \textbf{Der Invarianzkatalog ist eine Anfangsinventur}, keine
  Vollständigkeitsbehauptung. Insbesondere Arakelov- und adelische Grade bleiben
  unausgefüllt, weil das untersuchte Material nichts Verwertbares über sie sagt und
  keine Referenz geraten wurde.
\item \textbf{Textgrundlage und visuelle Rückprüfung.} Zitate verwenden
  \texttt{pdftotext}-Extrakte; Zeilenangaben sind Auffindhilfen, keine Paginierung.
  Die Quell-PDF-Seiten zu Gleichung (5.3.3), Theorem~5.10.1(3)--(4),
  Theorem~10.20.1(3), Definition~4.4.2/Formel (4.4.5) und
  Theorem~6.10.1 wurden visuell geprüft. Tragend sind insbesondere die
  Exponentenrichtung und beide Paare der E1-Betragsstriche. Für andere Passagen
  wird keine pauschale OCR-Robustheit behauptet.
\item \textbf{Keine Anhebung von Behauptungen.} Nichts in dieser Arbeit ist eine
  Aussage über die Korrektheit von Joshis $abc$-Argument, über Mochizukis Theorie
  \cite{Mochizuki2021a,Mochizuki2021b,Mochizuki2021c,Mochizuki2021d}, über die
  Einwände von~\cite{ScholzeStix2018} oder über die Riemannsche Vermutung. Sie ist
  die Klärung einer definitorischen Frage und ihrer Konsequenz für ein vermutetes
  Lemma.
\end{enumerate}

% =============================================================================
% KI-OFFENLEGUNG -- Pipeline-Standard, eingebunden aus der kanonischen Vorlage
% (_templates/disclosure/ai_disclosure_STANDARD.tex, zweisprachig, bringt eigenes
% \section* mit). Die Vorlage nennt Rechenskripte und Berechnungsläufe; diese
% Arbeit hat keine, daher passt der Einschränkungsabsatz darunter sie nach unten
% an -- genau das verlangt die Anweisung der Vorlage selbst ("Abweichen nach
% unten"). Die Einschränkung nicht stillschweigend weglassen: ohne sie behauptete
% die Offenlegung nicht geleistete Arbeit. Identisch in der englischen Fassung.
% =============================================================================
\phantomsection
\addcontentsline{toc}{section}{KI-Offenlegung}
% =============================================================================
% AI DISCLOSURE -- PIPELINE-STANDARD, DEUTSCHE FASSUNG (seit 2026-08-17, T-20260817-02)
% =============================================================================
% Sprachreine Fassung fuer rein deutschsprachige Manuskripte. Text 1:1 aus dem
% deutschen Block von ai_disclosure_STANDARD.tex uebernommen (keine inhaltliche
% Aenderung) -- nur die Ueberschrift ist jetzt sprachrein und der englische
% Block entfaellt. Die bilinguale Datei ai_disclosure_STANDARD.tex bleibt
% unveraendert bestehen (Bestandsverweise duerfen weiter darauf zeigen).
%
% ABWEICHEN NACH UNTEN: Wurde in einem Paper tatsaechlich deutlich weniger
% KI eingesetzt, wird NICHT auf eine andere Datei gewechselt, sondern die
% Formulierung hier abgeschwaecht -- z. B. "in erheblichem Masse in allen
% Phasen" ersetzen durch die tatsaechlich zutreffenden Phasen, und aus der
% Qualitaetssicherungs-Aufzaehlung streichen, was nicht stattfand.
% Der Grundsatz bleibt: ehrlich benennen, was war -- nicht abgrenzen.
%
% Verwendung (nur fuer deutschsprachige Manuskripte):
%   \input{../../_templates/disclosure/ai_disclosure_STANDARD_de}
% =============================================================================

\section*{KI-Offenlegung}

\noindent
KI-Systeme wurden in erheblichem Ma\ss{}e in allen Phasen dieser Arbeit
eingesetzt: Konzeption und Forschungsdiskussion, Literatur- und
Quellenarbeit, mathematische Analyse und Beweisarbeit, Erstellung und
Pr\"ufung der Rechenskripte, Ausf\"uhrung und Auswertung der
Berechnungsl\"aufe, Text, \"Ubersetzung und Satz.

\medskip

\noindent
Qualit\"atssicherung: Einsatz mehrerer unabh\"angiger Modelle
verschiedener Anbieter mit Gegenpr\"ufungen und adversarialen KI-Reviews;
nachvollziehbare, versionierte Rechenskripte mit gespeicherten
Ergebnisdateien und Pr\"ufsummen; fortlaufende Beweis- und
Registerdokumente; strikte Claim-Hygiene (Behauptungen nur mit
dokumentiertem Beleg, Nachtr\"age statt stiller Korrekturen); Changelogs
\"uber alle Versionen.


\medskip

\noindent
\emph{Einschränkung für diese Arbeit:} Zu dieser Arbeit gehören keine Rechenskripte
und keine Berechnungsläufe; die entsprechenden Punkte des Standardtextes
(Erstellung und Prüfung der Rechenskripte, Ausführung und Auswertung der
Berechnungsläufe, versionierte Skripte mit Ergebnisdateien und Prüfsummen) treffen
hier nicht zu.

% =============================================================================
% BIBLIOGRAPHY
% Every entry is provisional and must be checked against the original before any
% upload or submission. No entry was included that is not used in the text or
% expressly referenced by the working notes.
% =============================================================================
\begin{thebibliography}{99}
\addcontentsline{toc}{section}{Literatur}
\small

% verified 2026-08-13 — Quellencheck (WebSearch bzw. lokale Extrakte), Beleg: _sources/QUELLENCHECK_BIBITEMS_2026-08-13.md
\bibitem{JoshiATSIIhalf}
K.~Joshi.
\newblock \emph{Construction of Arithmetic Teichmüller Spaces II(1/2):
Deformations of Number Fields}.
\newblock arXiv:2305.10398, version v12 (24 February 2025). Hier zitiert als
ATS~II(1/2); \glqq [Joshi, 2023a]\grqq{} in den Literaturverzeichnissen des Autors.

% verified 2026-08-13 — Quellencheck (WebSearch bzw. lokale Extrakte), Beleg: _sources/QUELLENCHECK_BIBITEMS_2026-08-13.md
\bibitem{JoshiATSII}
K.~Joshi.
\newblock \emph{Construction of Arithmetic Teichmüller Spaces II: Proof of a
local prototype of Mochizuki's Corollary 3.12}.
\newblock arXiv:2303.01662, version v3 (24 February 2025). Hier zitiert als
ATS~II; \glqq [Joshi, 2023b]\grqq{} in den Literaturverzeichnissen des Autors.

% verified 2026-08-13 against local extract header (arXiv:2401.13508v4)
\bibitem{JoshiATSIII}
K.~Joshi.
\newblock \emph{Construction of Arithmetic Teichmüller Spaces III: A `Rosetta
Stone' and a proof of Mochizuki's Corollary 3.12}.
\newblock arXiv:2401.13508, version v4 (24 February 2025). Hier zitiert als
ATS~III; \glqq [Joshi, 2024c]\grqq{} in den Literaturverzeichnissen des Autors.

% verified 2026-08-13 against local extract header (arXiv:2403.10430v2)
\bibitem{JoshiATSIV}
K.~Joshi.
\newblock \emph{Construction of Arithmetic Teichmüller Spaces IV: Proof of the
abc-conjecture}. Preliminary version for comments.
\newblock arXiv:2403.10430, version v2 (24 February 2025). Hier zitiert als ATS~IV.

% verified 2026-08-13 — Quellencheck (WebSearch bzw. lokale Extrakte), Beleg: _sources/QUELLENCHECK_BIBITEMS_2026-08-13.md
\bibitem{Mochizuki2021a}
S.~Mochizuki.
\newblock \emph{Inter-universal Teichmüller Theory I: Construction of Hodge
Theaters}.
\newblock Publ. Res. Inst. Math. Sci. \textbf{57} (2021), no.~1/2, 3--207.

% verified 2026-08-13 — Quellencheck (WebSearch bzw. lokale Extrakte), Beleg: _sources/QUELLENCHECK_BIBITEMS_2026-08-13.md
\bibitem{Mochizuki2021b}
S.~Mochizuki.
\newblock \emph{Inter-universal Teichmüller Theory II: Hodge--Arakelov-theoretic
Evaluation}.
\newblock Publ. Res. Inst. Math. Sci. \textbf{57} (2021), no.~1/2, 209--401.

% verified 2026-08-13 — Quellencheck (WebSearch bzw. lokale Extrakte), Beleg: _sources/QUELLENCHECK_BIBITEMS_2026-08-13.md
\bibitem{Mochizuki2021c}
S.~Mochizuki.
\newblock \emph{Inter-universal Teichmüller Theory III: Canonical Splittings of
the Log-theta-lattice}.
\newblock Publ. Res. Inst. Math. Sci. \textbf{57} (2021), no.~1/2, 403--626.
\newblock Korollar~3.12 und Bemerkung~3.1.1 sind die über Joshis Text
referenzierten Passagen.

% verified 2026-08-13 — Quellencheck (WebSearch bzw. lokale Extrakte), Beleg: _sources/QUELLENCHECK_BIBITEMS_2026-08-13.md
\bibitem{Mochizuki2021d}
S.~Mochizuki.
\newblock \emph{Inter-universal Teichmüller Theory IV: Log-volume Computations
and Set-theoretic Foundations}.
\newblock Publ. Res. Inst. Math. Sci. \textbf{57} (2021), no.~1/2, 627--723.

% verified 2026-08-13 — Quellencheck (WebSearch bzw. lokale Extrakte), Beleg: _sources/QUELLENCHECK_BIBITEMS_2026-08-13.md
\bibitem{ScholzeStix2018}
P.~Scholze and J.~Stix.
\newblock \emph{Why abc is still a conjecture}.
\newblock Manuskript, 2018. Verfügbar unter
\texttt{https://www.math.uni-bonn.de/people/scholze/WhyABCisStillaConjecture.pdf}.

% verified 2026-08-13 — Quellencheck (WebSearch bzw. lokale Extrakte), Beleg: _sources/QUELLENCHECK_BIBITEMS_2026-08-13.md
\bibitem{ArtinWhaples1945}
E.~Artin and G.~Whaples.
\newblock \emph{Axiomatic characterization of fields by the product formula for
valuations}.
\newblock Bull. Amer. Math. Soc. \textbf{51} (1945), 469--492.

% v1.1: einziger neuer Eintrag, aufgenommen, weil Lemma 1 und Abschnitt 3.3(i)
% eine verifizierbare Fundstelle für die klassische Eindeutigkeitsaussage
% brauchen. Fundstelle [V]: PDF eingesehen, Prop. 7.8 auf S. 108 im Wortlaut
% geprüft (2026-08-13). Milnes Nummerierung ist versionsabhängig -- Version und
% Seite müssen mitzitiert bleiben.
\bibitem{Milne2020}
J.~S.~Milne.
\newblock \emph{Algebraic Number Theory}.
\newblock Skript, Version 3.08 (19. Juli 2020). Verfügbar unter
\texttt{https://www.jmilne.org/math/CourseNotes/ANT.pdf}. Zitiert mit Version,
Nummer und Seite, da die Nummerierung zwischen Versionen abweicht.

% verified 2026-08-13 — Quellencheck (WebSearch bzw. lokale Extrakte), Beleg: _sources/QUELLENCHECK_BIBITEMS_2026-08-13.md
\bibitem{FarguesFontaine2018}
L.~Fargues and J.-M.~Fontaine.
\newblock \emph{Courbes et fibr\'es vectoriels en th\'eorie de Hodge $p$-adique}.
\newblock Ast\'erisque \textbf{406}, Soci\'et\'e Math\'ematique de France, 2018.
Pr\'eface par Pierre Colmez, xiii+382~pp.

% verified 2026-08-13 — Quellencheck (WebSearch bzw. lokale Extrakte), Beleg: _sources/QUELLENCHECK_BIBITEMS_2026-08-13.md
\bibitem{BombieriGubler2006}
E.~Bombieri and W.~Gubler.
\newblock \emph{Heights in Diophantine Geometry}.
\newblock New Mathematical Monographs 4, Cambridge University Press, 2006.

% verified 2026-08-13 — Quellencheck (WebSearch bzw. lokale Extrakte), Beleg: _sources/QUELLENCHECK_BIBITEMS_2026-08-13.md
\bibitem{GranvilleStark2000}
A.~Granville and H.~M.~Stark.
\newblock \emph{ABC implies no ``Siegel zeros'' for $L$-functions of characters
with negative discriminant}.
\newblock Invent. Math. \textbf{139} (2000), 509--523.

\end{thebibliography}

\end{document}
